\documentclass[times,final]{elsarticle}
%\documentclass[times,twocolumn,final]{elsarticle}%
%\documentclass[times,twocolumn,final,longtitle]{elsarticle}
%%
%% Stylefile to load MEDIMA template
\usepackage{medima}
\usepackage{framed,multirow}
%% The amssymb package provides various useful mathematical symbols
\usepackage{amssymb}
\usepackage{amsmath}
\usepackage{latexsym}
\usepackage[left]{lineno}
\usepackage{subfigure}
\usepackage[utf8]{inputenc}
\usepackage[T1]{fontenc}
\usepackage[brazilian]{babel}
\usepackage{graphicx}
\usepackage{natbib}
\usepackage{booktabs}
\usepackage{float}
\usepackage{tikz}
\usetikzlibrary{calc,positioning,arrows.meta}


\graphicspath{{figuras/}}


\usepackage{url}
\usepackage{xcolor}
\usepackage{hyperref}
\definecolor{newcolor}{rgb}{.8,.349,.1}
\journal{Measurement}
\begin{document}
\verso{Silva, Calixto \textit{et~al.}}
\begin{frontmatter}
\title{Contribuição incremental das inter-harmônicas para diagnóstico eletromecânico e separação espectral de cargas em sistemas trifásicos}


\author[1,4,5]{Filipe Fraga Paula \snm{Silva}}
\ead{filipefragaps@gmail.com}
\author[1,4,5]{José A. Gobbes \snm{Cararo}}
\author[1,3,4]{Viviane M. \snm{Gomes Pacheco}}
\author[3]{Cloves Gonçalves \snm{Rodrigues}}
\author[2]{Antonio Paulo \snm{Coimbra}}
\author[1,2,5]{Wesley~Pacheco~\snm{Calixto}}
\ead{wesley.pacheco@ufg.br}

\address[1]{Electrical, Mechanical \& Computer Engineering School, Federal University of Goias, Goiania, Goias/Brazil. PC~74.605-010}
\address[2]{Institute of Systems and Robotics, Coimbra University, Coimbra/Portugal. PC~3030-290}
\address[3]{Polytechnic and Arts School, Pontifical Catholic University of Goias, Goiania, Goias/Brazil. PC~74.605-010}
\address[4]{Agroindustrial Automation and Precision Agriculture (AutoAgri), Federal Institute Goiano, Postal code: 75389-269, Trindade/Brazil}
\address[5]{Technology Research and Development Center, Federal Institute of Goias, Goiania, Goias/Brazil, PC~75.250-000}

\received{\today}
\finalform{xx xxx 2025}
\accepted{xx xxx 2025}
\availableonline{xx xxx 2025}
\communicated{S. Sarkar}

\begin{abstract}
\textcolor{blue}{Este trabalho investiga se uma mesma assinatura espectral agregada -- componentes harmônicas e inter-harmônicas de tensão e corrente, extraídas por um analisador de qualidade de energia já presente na instalação, sem forma de onda bruta nem sensoriamento dedicado -- é suficiente para responder, simultaneamente, ao diagnóstico do modo de falha de um sistema trifásico e à desagregação espectral de cargas elétricas coexistentes no mesmo ponto de acoplamento comum. Os dados foram coletados em duas bancadas complementares: um conjunto moto-máquina trifásico sob condições saudável, falha elétrica isolada, falha mecânica isolada (desbalanceamento e desalinhamento) e falha combinada, totalizando 17 condições experimentais e 19.501 aquisições; e um retificador trifásico alimentando combinações controladas de cargas não lineares. A partir de um espaço unificado de 984 canais espectrais candidatos, seis modelos clássicos de referência tabular (incluindo um Perceptron Multicamadas, tratado à parte por não possuir mecanismo de recorrência ou atenção) e seis arquiteturas sequenciais de aprendizado profundo -- incluindo uma arquitetura baseada em atenção -- foram comparados para o diagnóstico de falha sob um protocolo de validação temporal, desenhado para evitar vazamento de informação entre aquisições vizinhas da mesma sessão de coleta; o melhor modelo atinge 88--90\% de acurácia média sob validação cruzada em blocos temporais. Para a desagregação de cargas, um regressor \textit{Random Forest} multi-saída é empregado, dado o número reduzido de condições combinadas disponíveis para treino. Uma ablação controlada mostra diretamente que o subconjunto combinado harmônicas+inter-harmônicas supera o subconjunto apenas-harmônicas tanto no diagnóstico quanto na desagregação de cargas -- nesta última, avaliada sob seleção de atributos sem vazamento, o canal inter-harmônico mais recuperável atinge coeficiente de determinação de 0,88, evidência direta de que componentes inter-harmônicas, sistematicamente omitidas na literatura de monitoramento não intrusivo de cargas revisada, carregam informação diagnóstica relevante. A desagregação de cargas, avaliada com apenas quatro condições combinadas disponíveis, permanece um estudo de viabilidade parcial e seletiva por canal; e a evidência de previsão de falha um passo à frente não superou um preditor de persistência trivial, de modo que os resultados sustentam diagnóstico de condição e manutenção baseada em condição, não ainda detecção precoce de falha em desenvolvimento -- limitações registradas explicitamente, ao lado da ausência de defeitos de rolamento na base de dados atual, como direções de trabalho futuro.}
\end{abstract}
\begin{keyword}
\textcolor{blue}{Inter-harmônicas \sep diagnóstico de falhas em máquinas rotativas \sep monitoramento não intrusivo de cargas (NILM) \sep análise de assinatura de corrente do motor (MCSA) \sep aprendizado profundo \sep qualidade de energia.}
\end{keyword}
\end{frontmatter}
\linenumbers
%===============================================================
\section{Introduction}
\label{cap:introducao}
O crescimento contínuo da eletrônica de potência e da automação industrial tem alterado de forma expressiva o comportamento espectral das grandezas elétricas em sistemas de distribuição de baixa e média tensão. Conversores estáticos, retificadores, acionamentos de velocidade variável e demais cargas não lineares injetam, além da componente fundamental de 50/60~Hz, um espectro de componentes harmônicas -- múltiplos inteiros da frequência fundamental -- e, em determinadas condições operacionais, também componentes interharmônicas, definidas pela norma IEC~61000-4-7 como aquelas cuja frequência não corresponde a um múltiplo inteiro da fundamental \citep{efe2016}. Esse cenário, tratado por décadas majoritariamente como um problema de qualidade de energia a ser mitigado, tem sido crescentemente reinterpretado pela literatura como uma fonte de informação: a distorção harmônica e interharmônica carrega, em sua estrutura, assinaturas que podem ser exploradas tanto para identificar cargas elétricas individuais operando em um mesmo ponto de acoplamento comum quanto, por extensão, para inferir condições de operação de equipamentos rotativos conectados ao sistema.

\textcolor{blue}{
\subsection{Problema, motivação e lacuna}
Apesar dessa reinterpretação do conteúdo espectral como fonte de informação, dois problemas práticos permanecem em aberto na indústria: identificar, a partir de um único ponto de medição já instalado, quais cargas elétricas não lineares compõem o consumo agregado de um sistema trifásico; e diagnosticar a condição de máquinas rotativas conectadas a esse mesmo sistema, idealmente antes que uma falha em desenvolvimento evolua para uma parada não planejada. A literatura de monitoramento não intrusivo de cargas (NILM) baseado em conteúdo harmônico e a literatura clássica de diagnóstico de falha via análise de assinatura de corrente (MCSA) -- revisadas em detalhe no Capítulo~\ref{cap:fundamentacao} -- endereçam, cada uma isoladamente, um desses dois problemas a partir de princípio semelhante: extrair, do conteúdo espectral da corrente (e, em menor medida, da tensão) já mensurável em pontos de medição existentes, uma assinatura suficiente para identificação ou diagnóstico, sem exigir sensoriamento dedicado adicional. A motivação central desta pesquisa é que essa mesma premissa de não intrusividade permite, em princípio, tratar diagnóstico de falha e desagregação de cargas como duas saídas de um único espaço de entrada espectral agregado -- possibilidade que a leitura crítica da literatura (Seção~\ref{sec:sintese}) mostra não ter sido testada por nenhum dos trabalhos revisados.
Essa leitura crítica revela ainda uma segunda lacuna, ortogonal à primeira: nenhum dos doze trabalhos de NILM e de MCSA revisados no Capítulo~\ref{cap:fundamentacao} avalia sistematicamente o conteúdo interharmônico -- componentes cuja frequência não corresponde a um múltiplo inteiro da fundamental \citep{efe2016} -- como atributo de identificação de carga ou de diagnóstico de falha, apesar de \citet{efe2016} demonstrar, em um sistema retificador estruturalmente análogo ao desta pesquisa, que interharmônicas podem apresentar distorção superior à das harmônicas adjacentes. É a combinação dessas duas lacunas -- a ausência de avaliação sistemática de componentes interharmônicas, e a ausência de um sistema unificado de diagnóstico e desagregação em ambiente trifásico industrial -- que motiva o problema de pesquisa e a hipótese formulados a seguir.

\subsection{Hipótese, objetivo geral e objetivos específicos}
\label{sec:hipotese}
A hipótese central desta pesquisa, testada diretamente por ablação controlada (Seção~\ref{sec:resultados_ablacao}), é:
\begin{quote}
$H_1$: $M_{H+IH} > M_H$ -- o desempenho de um modelo treinado sobre o espaço de atributos combinado (harmônicas $+$ inter-harmônicas) é superior ao de um modelo treinado apenas sobre harmônicas, sob um mesmo protocolo de validação pré-definido (Seção~\ref{sec:protocolo}), tanto para o diagnóstico de falha (RQ1, $M=$ F1-macro) quanto para a desagregação de cargas (RQ2, $M=R^2$ por canal).
\end{quote}
A partir dessa hipótese, formula-se o seguinte objetivo geral: \textbf{desenvolver e avaliar uma abordagem baseada em assinaturas espectrais agregadas (harmônicas e inter-harmônicas de tensão e corrente) para diagnóstico de condições eletromecânicas e separação de contribuições espectrais em sistemas trifásicos}, sem recorrer a sensoriamento dedicado ou à forma de onda bruta. Esse objetivo geral se desdobra nos seguintes objetivos específicos, cada um vinculado a um experimento verificável desenvolvido nesta pesquisa:
\begin{itemize}
    \item OE1: construir um espaço de atributos espectrais unificado (984 canais candidatos, harmônicos e inter-harmônicos, de tensão e corrente) comum aos dois problemas de pesquisa (Seção~\ref{sec:atributos});
    \item OE2: comparar sistematicamente modelos clássicos de aprendizado de máquina e arquiteturas sequenciais de aprendizado profundo para o diagnóstico de falha (RQ1), sob um protocolo de validação temporal sem vazamento de informação entre treino e teste (Seções~\ref{sec:protocolo} e~\ref{sec:resultados_rq1});
    \item OE3: avaliar a viabilidade de desagregar espectralmente as contribuições de cargas elétricas coexistentes no mesmo ponto de acoplamento comum (RQ2), a partir exclusivamente da mesma assinatura agregada (Seção~\ref{sec:resultados_rq2});
    \item OE4: testar diretamente a hipótese $H_1$ por ablação controlada -- comparando, sob dobras idênticas, subconjuntos de atributos apenas-harmônicos, apenas-interharmônicos e combinados -- tanto para RQ1 quanto para RQ2 (Seção~\ref{sec:resultados_ablacao});
    \item OE5: interpretar os modelos treinados via SHAP, identificando quais ordens harmônicas e interharmônicas, de tensão e de corrente, mais contribuem para cada tarefa (Seção~\ref{sec:interpretabilidade}).
\end{itemize}
Esses objetivos específicos se organizam em torno das duas perguntas de pesquisa co-primárias que fundamentam esta tese, detalhadas e contextualizadas frente à literatura revisada na Seção~\ref{sec:objetivos} do Capítulo~\ref{cap:fundamentacao}:
\textbf{RQ1 (diagnóstico e manutenção preditiva)}: um modelo treinado sobre a assinatura espectral agregada consegue diagnosticar corretamente o modo de falha de um sistema trifásico -- saudável, falha elétrica isolada, falha mecânica isolada ou falha combinada -- com desempenho suficiente para apoiar decisões de manutenção baseada em condição?
\textbf{RQ2 (desagregação de cargas)}: com cargas elétricas operando simultaneamente no mesmo ponto de acoplamento comum, é possível desagregar suas contribuições espectrais individuais a partir exclusivamente da assinatura agregada, e as componentes interharmônicas contribuem de forma relevante para essa desagregação?

\subsection{Contribuições pretendidas e escopo}
Frente às lacunas identificadas, esta pesquisa pretende contribuir de três formas, retomadas e avaliadas criticamente à luz dos resultados obtidos na Seção~\ref{sec:discussao_contribuicoes}: (i) tratar diagnóstico de falha mecânica e desagregação de cargas elétricas como duas saídas de um mesmo espaço de entrada espectral agregado, ao invés de dois sistemas independentes; (ii) transformar a relevância diagnóstica das interharmônicas -- até então sugerida apenas indiretamente na literatura \citep{efe2016} -- em uma hipótese formal ($H_1$) diretamente testável por ablação controlada, em vez de inferida apenas de rankings de importância de atributos; e (iii) documentar, com o mesmo rigor metodológico exigido de um problema de pesquisa autônomo, os riscos de validação específicos de dados de séries temporais industriais com sessões de aquisição únicas por condição (Seção~\ref{sec:protocolo}). O escopo experimental está delimitado a duas bancadas complementares -- diagnóstico de falha em motor de indução trifásico e desagregação em retificador trifásico (Seção~\ref{sec:bancadas}) --, totalizando 17 condições experimentais e 19.501 aquisições (Seção~\ref{sec:resultados_bd}); a delimitação exata desse escopo, incluindo o que fica deliberadamente fora dele (por exemplo, defeitos de rolamento e replicação em outras instalações), é discutida na Seção~\ref{cap:limitacoes}.

\subsection{Estrutura do documento}
O restante deste documento está organizado da seguinte forma: o Capítulo~\ref{cap:fundamentacao} revisa a literatura de NILM, MCSA e os fundamentos de aprendizado de máquina/profundo empregados; o Capítulo~\ref{cap:metodologia} descreve a metodologia experimental, a base de dados e o protocolo de análise estatística e de aprendizado de máquina utilizados; o Capítulo~\ref{cap:resultados} apresenta os resultados obtidos para RQ1 e RQ2, incluindo a ablação controlada que testa $H_1$; o Capítulo~\ref{cap:discussao} discute esses resultados à luz da literatura revisada e registra as limitações estruturais da pesquisa; o Capítulo~\ref{cap:conclusao} sintetiza as respostas a RQ1 e RQ2 e conclui o documento; e o Apêndice~B consolida um checklist de reprodutibilidade dos experimentos.
}
%===============================================================
\section{Theoretical Background}
\label{cap:fundamentacao}
Este capítulo revisa criticamente duas vertentes de literatura necessárias à fundamentação da pesquisa: o monitoramento não intrusivo de cargas (NILM) baseado em conteúdo harmônico, e o diagnóstico de falhas em máquinas rotativas via análise de assinatura de corrente (MCSA). As duas vertentes convergem, como discutido na Seção~\ref{sec:sintese}, para uma lacuna comum -- a ausência de avaliação sistemática de componentes interharmônicas como atributo diagnóstico -- que motiva as perguntas de pesquisa formalizadas na Seção~\ref{sec:objetivos}, ao final deste capítulo. Complementarmente, a Seção~\ref{sec:ml_dl} revisa os fundamentos conceituais e históricos dos métodos de aprendizado de máquina e aprendizado profundo empregados na metodologia experimental (Capítulo~\ref{cap:metodologia}), estabelecendo a base técnica sobre a qual os resultados de diagnóstico e desagregação (Capítulo~\ref{cap:resultados}) são posteriormente discutidos.
\subsection{Monitoramento não intrusivo de cargas baseado em conteúdo harmônico}
O paradigma de monitoramento não intrusivo de cargas (\textit{Non-Intrusive Load Monitoring} -- NILM) parte da premissa de que é possível desagregar o consumo agregado, medido em um único ponto de medição, nas contribuições individuais de cada carga ativa, sem necessidade de sensoriamento distribuído. O campo tem origem no trabalho seminal de \citet{hart1992}, que propôs originalmente a identificação de eletrodomésticos a partir de transições em degrau de potência ativa e reativa observadas em um único ponto de medição, estabelecendo o paradigma não intrusivo como alternativa de baixo custo ao sensoriamento distribuído então predominante. Quase duas décadas depois, \citet{zeifman2011} revisam a evolução do campo a partir dessa proposta original, sistematizando as principais famílias de técnicas de desagregação desenvolvidas no intervalo -- baseadas em regime permanente, como a variação em degrau de potência, e em regime transitório -- e apontando, já naquele momento, a necessidade de atributos mais ricos do que a simples variação de potência ativa e reativa para discriminar cargas de comportamento espectral semelhante, apontamento que antecipa a linha de investigação retomada por \citet{djordjevic2017}. Conforme discutido por \citet{djordjevic2017}, o método clássico de identificação de cargas baseia-se nas variações em degrau de potência ativa e reativa, mapeadas em um plano bidimensional P-Q; essa abordagem, no entanto, apresenta desempenho limitado para cargas não lineares de baixa potência, cujas assinaturas tendem a se sobrepor no plano P-Q, sobretudo quando o consumo é inferior a poucas centenas de watts.
Para superar essa limitação, \citet{djordjevic2017} propõem substituir -- ou complementar -- as grandezas de potência por uma única grandeza de corrente harmônica agregada ($I_H$), calculada a partir da diferença entre o valor eficaz total da corrente e a componente fundamental, sem exigir o cálculo individual de cada ordem harmônica. Os autores demonstram que a separação de assinaturas no plano potência ativa--corrente harmônica ($P$--$I_H$) é substancialmente mais nítida do que no plano $P$--$Q$ para o mesmo conjunto de cargas residenciais não lineares, inclusive para eletrodomésticos de potência semelhante, e que o custo computacional dessa abordagem é inferior ao de métodos baseados no cálculo de múltiplas ordens harmônicas via FFT, uma vez que os autores recorrem ao algoritmo de Goertzel para extrair apenas a componente fundamental necessária ao cálculo de $I_H$.
\citet{bouhouras2017} avançam essa linha de investigação ao mostrar que uma única grandeza de corrente harmônica agregada nem sempre é suficiente para discriminar cargas com comportamento espectral mais complexo, propondo a formulação de assinaturas de carga (\textit{load signatures}) a partir de valores representativos das correntes na frequência fundamental e nas ordens harmônicas ímpares de 3\textsuperscript{a} e 5\textsuperscript{a} ordem, obtidos por um procedimento estatístico que agrupa os valores medidos conforme seu desvio-padrão. Os resultados experimentais, obtidos para combinações de até cinco cargas residenciais operando simultaneamente, indicam desempenho de identificação próximo de 100\% quando as cargas envolvidas possuem conteúdo harmônico pouco significativo, mas revelam uma limitação estrutural relevante: como as correntes harmônicas não se combinam linearmente em amplitude quando múltiplas cargas com conteúdo harmônico expressivo operam simultaneamente -- a soma vetorial depende da fase de cada componente, não considerada no modelo --, a taxa de identificação cai de forma acentuada nesses casos, chegando a 0\% para determinadas combinações. Esse achado é metodologicamente relevante para qualquer sistema, como o investigado nesta tese, em que a desagregação é realizada exclusivamente a partir de amplitudes espectrais agregadas, sem acesso à fase relativa entre as componentes de diferentes fontes.
\citet{li2019moea} deslocam o problema de desagregação para uma formulação multiobjetivo, integrando simultaneamente potência ativa, reativa, aparente, forma de onda de corrente e conteúdo harmônico como funções-objetivo distintas, resolvidas por algoritmos evolutivos multiobjetivo -- com destaque para o NSGA-II, que superou estatisticamente as demais variantes testadas (MOEA/D-MTCH, MOEA/D-PBI, MOEA/D-TCH, MOEA/D-TCHN e RVEA). O resultado central desse trabalho -- de que o aumento do número de assinaturas de carga utilizadas simultaneamente melhora consistentemente a taxa de detecção de soluções reais e reduz o erro de desagregação, com ganhos de até 83\% em relação a métodos de otimização com atributo único -- reforça, por uma via distinta da de \citet{bouhouras2017}, a hipótese de que a riqueza do espaço de atributos espectrais é o fator determinante do desempenho de desagregação, mais do que a escolha isolada de um único descritor ótimo.
\subsection{Interharmônicas como conteúdo espectral subexplorado}
Os três trabalhos revisados na seção anterior restringem suas assinaturas de carga a ordens harmônicas inteiras, sem considerar explicitamente componentes interharmônicas. Essa escolha metodológica é compreensível do ponto de vista de simplicidade computacional, mas carece de justificativa empírica direta quanto à real irrelevância diagnóstica dessas componentes. Nesse ponto, o trabalho de \citet{efe2016} oferece uma contribuição importante, ainda que não formulada originalmente para fins de identificação de cargas ou diagnóstico de falhas. Ao simular, em ambiente MATLAB/Simulink, um sistema composto por um retificador trifásico alimentando uma carga em corrente contínua -- topologia estruturalmente análoga à bancada de retificador utilizada na presente pesquisa --, o autor demonstra que componentes interharmônicas situadas nas proximidades de ordens harmônicas inteiras podem apresentar distorção percentual superior à da própria harmônica adjacente; por exemplo, a interharmônica de ordem 4,5 apresentou distorção de 1,27\%, superior à distorção de 0,34\% observada na 5\textsuperscript{a} harmônica no mesmo sistema. Esse resultado indica que a exclusão sistemática das interharmônicas nas formulações de assinatura de carga revisadas na Seção anterior pode representar o descarte de informação espectral potencialmente discriminativa -- hipótese que a presente pesquisa busca testar diretamente, ao incluir explicitamente componentes interharmônicas de tensão e corrente, lado a lado com as harmônicas, no espaço de atributos utilizado tanto para diagnóstico de falhas quanto para desagregação de cargas (RQ1 e RQ2, respectivamente, conforme detalhado na Seção~\ref{sec:objetivos}).
\subsection{Diagnóstico de falhas em máquinas rotativas via análise de assinatura de corrente (MCSA)}
A segunda vertente teórica necessária a esta tese diz respeito ao diagnóstico de condição de máquinas rotativas, historicamente consolidado sob a denominação de \textit{Motor Current Signature Analysis} (MCSA). A revisão de \citet{nandi2005} -- um dos trabalhos de maior impacto e mais citados nessa área -- estabelece a premissa central sobre a qual a técnica se apoia: diferentes modos de falha de uma máquina elétrica rotativa imprimem assinaturas fisicamente rastreáveis nas grandezas eletromagnéticas da máquina, sobretudo na corrente de estator. \textcolor{red}{Os autores reportam, com base em levantamentos setoriais anteriores, que falhas de rolamento respondem por aproximadamente 40\%--50\% de todas as falhas em motores de indução, falhas de isolamento/enrolamento estatórico por 30\%--40\% e falhas de rotor por 5\%--10\%, evidenciando que os modos de falha cobertos pela taxonomia experimental desta tese (desbalanceamento, desalinhamento, falha de rolamento e falha elétrica) correspondem, em conjunto, à quase totalidade dos modos de falha reportados na literatura de manutenção de motores. Os autores sistematizam ainda, de forma explícita, as relações de frequência que fundamentam a MCSA clássica: para quebra de barras do rotor, $f_b = (1 \pm 2s)f$ e, de forma mais geral, $f_b = [(k/p)(1-s) \pm s]f$; para excentricidade de entreferro, $f_e = f[(kR \pm n_d)(1-s)/p \pm \nu]$, em que $R$ é o número de ranhuras do rotor, $n_d$ o número de ordem da excentricidade (estática ou dinâmica), $p$ o número de pares de polos e $\nu$ a ordem das harmônicas de tensão presentes na alimentação; para falta entre espiras do enrolamento estatórico, $k(R/p) \pm 1$; e para falhas de rolamento, a tradução das frequências características de vibração ($f_\nu$, associadas às dimensões geométricas do rolamento e à frequência de rotação) para o domínio da corrente por meio de $f_{bng} = |f_1 \pm m \cdot f_\nu|$. Essas quatro famílias de equações constituem o repertório clássico de frequências-alvo da MCSA e serão retomadas, na Seção~\ref{sec:sintese}, como referência para contrastar a abordagem baseada em bandas laterais discretas com o espaço amplo de atributos espectrais agregados adotado nesta tese. É relevante notar, contudo, que os próprios autores reconhecem um limite estrutural dessa abordagem: ao tratar da falta entre espiras, afirmam explicitamente que ``a física por trás da existência de tais componentes não é claramente explicada'', e, em suas conclusões, ressaltam que ``a análise e a modelagem teóricas das falhas da máquina são de fato necessárias para distinguir os componentes de frequência relevantes de outros que podem estar presentes devido a harmônicas de tempo, saturação da máquina, etc.''. Essa ressalva é diretamente pertinente ao escopo desta tese: na ausência de um modelo físico-analítico prévio que isole univocamente cada componente de falha, a abordagem orientada a dados aqui adotada -- que trata o espectro agregado como espaço de atributos a ser explorado por classificadores, em vez de buscar isolar analiticamente bandas laterais específicas -- oferece um caminho alternativo, e não excludente, para lidar com essa mesma dificuldade.} Essa taxonomia é o ponto de partida conceitual para o problema de diagnóstico de falha (RQ1) desta tese, que amplia o alvo diagnóstico de falhas puramente elétricas para um espectro que inclui também falhas mecânicas isoladas e combinadas.
\citet{thomson2001} consolidam a MCSA como técnica não invasiva de diagnóstico já adotada pela indústria, afirmando que \textcolor{red}{``é atualmente aceito pela indústria que a MCSA é a técnica mais aplicável e eficaz para o diagnóstico de barras de rotor quebradas em motores de indução trifásicos''}, e demonstrando que a simples medição da corrente de estator -- tipicamente já disponível em centros de controle de motores ou em quadros de distribuição, sem necessidade de sensoriamento adicional junto à máquina -- é suficiente para detectar componentes de frequência características de quebra de barras do rotor e de excentricidade de entreferro, manifestas como bandas laterais (\textit{sidebands}) ao redor da frequência fundamental e de suas harmônicas\textcolor{red}{, segundo as relações $f_{sb} = f_1(1 \pm 2s)$ para quebra de barras, $f_{st} = f_1[(n/p)(1-s) \pm k]$ para faltas entre espiras e $f_{ec} = f_1[(R \pm n_d)(1-s)/p \pm n\omega_s]$ para excentricidade}. \textcolor{red}{Os autores sustentam essas afirmações com quatro estudos de caso industriais reais que ilustram tanto o poder diagnóstico quanto os riscos de má interpretação da técnica. No primeiro caso, dois motores idênticos de 134~kW acionando compressores foram comparados: um deles não apresentou bandas laterais (condição saudável), enquanto o outro exibiu bandas laterais a $\pm 0{,}72$~Hz da fundamental, com magnitude entre 32,5~dB e 36,2~dB abaixo da componente fundamental, levando à previsão de cinco a seis barras quebradas -- confirmada, após desmontagem do motor, pela constatação de seis a sete barras efetivamente fraturadas. No segundo caso, uma bomba de injeção de água do mar de 3,6~MW em uma plataforma offshore, inicialmente saudável (bandas laterais mais de 45~dB abaixo da fundamental), apresentou, após um incidente operacional, aumento de 26~dB na magnitude das bandas laterais -- fator de amplificação de vinte vezes em amplitude --, o que motivou a retirada preventiva do motor de serviço e a posterior confirmação de vinte barras quebradas em um rotor de 46 barras. No terceiro caso, um pulverizador de carvão de uma usina termelétrica apresentou componentes espectrais a $\pm 0{,}625$~Hz da fundamental, 33~dB abaixo desta -- padrão que, à primeira vista, se assemelharia a bandas laterais de quebra de barras, mas que os autores atribuem corretamente a componentes mecânicas do trem de acionamento (engrenagens), e não a uma falha de barras, ilustrando um risco de diagnóstico equivocado inerente à interpretação de bandas laterais isoladas sem contextualização do sistema mecânico acoplado. No quarto caso, um motor de 11~kW com falta entre espiras no enrolamento estatórico apresentou componentes a 125~Hz e 175~Hz, com magnitudes de 56,6~dB e 42,3~dB, respectivamente, permanecendo em operação por vinte minutos antes da falha final -- evidenciando uma janela de detecção precoce potencialmente explorável para prevenção de falhas catastróficas.} Esse resultado é diretamente relevante para a premissa metodológica desta tese, articulada já na contextualização do Capítulo~\ref{cap:introducao}: a de que a assinatura espectral agregada, obtida do mesmo analisador de energia já presente em instalações industriais, pode carregar informação diagnóstica sem instrumentação adicional. Cabe notar, contudo, uma divergência metodológica relevante: a MCSA clássica, tal como formulada por \citet{thomson2001}, concentra-se na detecção de bandas laterais estreitas associadas a mecanismos de falha elétrica específicos (relacionadas ao escorregamento da máquina)\textcolor{red}{ -- e, como o próprio terceiro estudo de caso demonstra, sujeita a ambiguidade de interpretação quando componentes mecânicas do sistema acoplado ocupam faixas de frequência próximas às bandas laterais elétricas --}, e não em um espaço amplo de atributos harmônicos e interharmônicos multifásicos tratado por classificadores de aprendizado de máquina, como o empregado nesta pesquisa.
\textcolor{red}{Uma atualização mais recente e diretamente relevante para o posicionamento desta tese é oferecida por \citet{niu2023}, cuja revisão sistemática sobre diagnóstico de falhas de motores baseado em assinaturas de corrente, publicada quase duas décadas após \citet{nandi2005}, sintetiza 142 trabalhos publicados na área e reflete a maturação e a diversificação metodológica do campo. Os autores organizam a literatura de aprendizado profundo aplicada à MCSA em quatro grandes famílias arquiteturais -- redes neurais convolucionais (CNN), autoencoders (AE) e suas variantes, redes de crença profunda (DBN) e redes adversariais generativas (GAN) --, notavelmente sem incluir, em nenhuma das quatro categorias, arquiteturas recorrentes (RNN, LSTM, GRU, BiLSTM) ou baseadas em atenção (Transformer), o que evidencia uma lacuna metodológica na literatura de MCSA especificamente -- em contraste com o diagnóstico de falhas de rolamento por vibração, no qual \citet{hoang2019} já identificavam essa mesma lacuna alguns anos antes. Essa ausência reforça diretamente a escolha desta tese de comparar sistematicamente seis arquiteturas sequenciais (SimpleRNN, LSTM, GRU, BiLSTM, CNN+LSTM e um Transformer temporal) para o problema de diagnóstico de falha (RQ1), preenchendo uma lacuna que persiste mesmo na revisão mais recente disponível sobre o tema. Adicionalmente, \citet{niu2023} não emprega em nenhum momento o termo ``interharmônica'' para descrever os componentes espectrais associados a falhas -- optando, como a quase totalidade da literatura de MCSA revisada nesta tese, pela terminologia de ``bandas laterais'' ou ``componentes de frequência característicos'' --, o que reforça o argumento, desenvolvido na Seção~\ref{sec:sintese}, de que a caracterização explícita de componentes interharmônicos no diagnóstico de falhas permanece uma lacuna terminológica e metodológica não endereçada mesmo pela revisão mais atual do campo. Os autores discutem ainda, como direção de pesquisa em aberto, o problema do acoplamento eletromecânico entre componentes de um sistema rotativo (por exemplo, motor e caixa de engrenagens, ou motor e carga mecânica): quando falhas ocorrem simultaneamente em componentes acoplados, as harmônicas de falha resultantes não se combinam de forma aditiva ou superponível, exigindo modelos de acoplamento dedicados em vez da simples sobreposição das assinaturas espectrais de cada falha isolada. Esse ponto é diretamente relevante para a categoria de falha combinada (mecânica + elétrica) desta tese, na medida em que alerta para o risco de que uma assinatura espectral combinada não seja uma simples superposição linear das assinaturas isoladas -- risco que a presente pesquisa mitiga ao tratar cada condição combinada como uma classe própria a ser aprendida diretamente pelos classificadores, em vez de inferida pela composição de assinaturas de falhas isoladas. Por fim, os autores apontam o caráter de caixa-preta das arquiteturas de aprendizado profundo como uma limitação relevante à adoção industrial dessas técnicas, argumento que converge com a motivação desta tese para a incorporação de métodos de interpretabilidade (SHAP), descrita no Capítulo~\ref{cap:metodologia}.}
Uma contribuição mais específica e diretamente relacionada ao escopo experimental desta tese é apresentada por \citet{benamira2022}, que investigam experimentalmente a ocorrência simultânea de desbalanceamento mecânico de massa e desequilíbrio de tensão de alimentação em um motor de indução trifásico, monitorando exclusivamente a corrente de estator. \textcolor{red}{O desbalanceamento mecânico é induzido pela fixação de uma massa adicional de 90~g ao rotor, enquanto o desequilíbrio de tensão é imposto eletricamente entre as três fases, com o fator de desequilíbrio de tensão (CUF) variando de 2,6\% (condição próxima do equilíbrio) até 21,17\% (condição de desequilíbrio severo) -- cabe notar que essa indução do desbalanceamento mecânico por fixação de massa adicional ao rotor difere metodologicamente da bancada empregada na presente pesquisa, na qual o desbalanceamento é obtido pela remoção controlada de parafusos de um disco calibrado, e não pela adição de massa.} Os autores demonstram, por meio de análise espectral das componentes direta e de quadratura do vetor de Park, que é possível não apenas detectar a ocorrência conjunta dessas duas falhas de naturezas distintas, mas também identificar qual fase específica apresentou a redução de tensão que originou o desequilíbrio elétrico\textcolor{red}{, associando a condição combinada a componentes espectrais específicas em torno de 25,44~Hz, 74,36~Hz e 149,7~Hz -- frequências que, por não corresponderem a múltiplos inteiros da fundamental de 50~Hz nem de suas harmônicas, seriam tecnicamente classificáveis como componentes interharmônicas segundo a terminologia de qualidade de energia (IEC/IEEE), ainda que os autores as descrevam exclusivamente como ``bandas laterais'' (\textit{sidebands}), sem nunca empregar o termo ``interharmônica'' -- reforçando, mais uma vez, a lacuna terminológica identificada nesta tese entre a literatura clássica de MCSA e a literatura de qualidade de energia}. Esse resultado oferece evidência direta, em um cenário experimental estruturalmente próximo ao desta tese, de que a categoria de falha combinada (mecânica + elétrica) adotada na taxonomia experimental do presente projeto é fisicamente plausível e espectralmente discriminável -- ainda que \citet{benamira2022} tratem de um único cenário combinado específico, e não de uma taxonomia com múltiplas condições combinadas nem de uma caracterização formal de componentes interharmônicas, que permanecem, nesse sentido, fora do escopo desse trabalho.
Complementando esse cenário, \citet{martinez2022} investigam experimentalmente a detecção de múltiplas falhas em motores de indução \textcolor{red}{-- especificamente cinco condições mutuamente isoladas: máquina saudável, uma barra de rotor quebrada, duas barras de rotor quebradas, falha de rolamento e desbalanceamento mecânico de massa, sem incluir uma condição de falha combinada no conjunto de dados experimentais --} a partir de indicadores estatísticos extraídos do regime transitório de partida da corrente de estator, especificamente a homogeneidade \textcolor{red}{(derivada da matriz de coocorrência em níveis de cinza -- GLCM -- aplicada ao sinal de corrente)} e a curtose \textcolor{red}{(momento estatístico de quarta ordem)} do sinal, sem recorrer à análise espectral em regime permanente característica da MCSA clássica. Os autores demonstram, \textcolor{red}{empregando um classificador do tipo perceptron multicamadas (MLP) sobre esse par de atributos,} que esses indicadores estatísticos no domínio do tempo permitem distinguir, \textcolor{red}{com acurácia de 100\% entre as cinco condições avaliadas}, a partir de um único sinal de corrente monitorado de forma não invasiva. Esse resultado é particularmente relevante para a presente tese porque oferece evidência direta e recente de viabilidade da detecção de desbalanceamento mecânico isolado \textcolor{red}{a partir exclusivamente da corrente de estator, por uma via metodológica distinta -- indicadores estatísticos do regime transitório, e não conteúdo espectral em regime permanente -- da adotada nesta pesquisa; a convergência de achados entre abordagens metodologicamente distintas (a de \citet{martinez2022}, restrita à falha isolada, e a de \citet{benamira2022}, que trata do desbalanceamento em cenário combinado com desequilíbrio de tensão) reforça a plausibilidade física da hipótese de que o desbalanceamento mecânico deixa assinatura identificável na corrente elétrica, tanto isolada quanto combinada a outras falhas}, ainda que a caracterização dessa assinatura no domínio espectral -- incluindo eventuais componentes interharmônicas -- permaneça, como discutido na Seção~\ref{sec:sintese}, uma lacuna específica a ser abordada por esta tese.
Por fim, no que se refere aos defeitos em rolamentos -- previstos no desenho experimental do projeto, mas ainda não incorporados à base de dados analisada --, \citet{hoang2019} revisam sistematicamente a aplicação de arquiteturas de aprendizado profundo (autoencoders e suas variantes empilhadas, máquinas de Boltzmann restritas/redes de crença profunda e redes neurais convolucionais) ao diagnóstico de falha de rolamentos, a partir de sinais de vibração, emissão acústica e, em menor medida, corrente do motor. Os autores argumentam que essas arquiteturas superam a dependência de extração manual de atributos que caracteriza os métodos clássicos de aprendizado de máquina, aprendendo representações hierárquicas diretamente dos dados brutos ou de representações de baixo nível extraídas nos domínios do tempo, da frequência ou tempo-frequência. \textcolor{red}{A revisão apoia-se predominantemente em conjuntos de dados públicos de referência -- notadamente os \textit{datasets} de rolamento da Case Western Reserve University (CWRU) e da Universidade de Cincinnati --, ambos estruturados para representar modos de falha isolados, sem condições de falha combinada.} Dois achados dessa revisão são particularmente relevantes para o posicionamento metodológico desta tese. Primeiro, os autores observam que, apesar de os sinais de rolamento serem inerentemente séries temporais, havia até a publicação da revisão uma escassez notável de trabalhos aplicando redes neurais recorrentes (RNN) ao diagnóstico de falha de rolamentos\textcolor{red}{ -- lacuna mencionada pelos autores apenas como direção futura de pesquisa, na conclusão do artigo, sem que nenhuma arquitetura recorrente tenha sido efetivamente revisada ao longo do trabalho --}, apontando-as como ferramenta promissora e ainda pouco explorada nesse domínio específico -- lacuna que a presente tese endereça diretamente ao empregar e comparar sistematicamente seis arquiteturas sequenciais (SimpleRNN, LSTM, GRU, BiLSTM, CNN+LSTM e um Transformer temporal) para o diagnóstico de falha (RQ1). Segundo, os autores apontam que a quase totalidade dos trabalhos revisados trata exclusivamente do diagnóstico de falha única\textcolor{red}{ -- consistente com a própria natureza dos conjuntos de dados de referência empregados (CWRU, Cincinnati), estruturalmente limitados a condições de falha isolada --}, permanecendo em aberto a verificação da viabilidade do aprendizado profundo em cenários de múltiplas falhas simultâneas -- desafio que a taxonomia experimental desta tese já incorpora ao tratar explicitamente de uma categoria de falha combinada (mecânica + elétrica), ainda que, como já mencionado, defeitos em rolamentos especificamente não integrem, por ora, a base de dados analisada.
\subsection{Fundamentos dos métodos de aprendizado de máquina e aprendizado profundo empregados nesta pesquisa}
\label{sec:ml_dl}
Além da fundamentação específica sobre NILM e MCSA apresentada nas seções anteriores, a metodologia desta tese apoia-se em um conjunto de algoritmos de aprendizado de máquina e aprendizado profundo cuja origem conceitual e cujo desenvolvimento histórico são brevemente revisados nesta seção, na medida em que constituem a base técnica direta dos modelos comparados nos Capítulos~\ref{cap:metodologia} e~\ref{cap:resultados}.
\citet{lecun2015} sintetizam, em um dos artigos de maior impacto da área, os fundamentos e a trajetória do aprendizado profundo, caracterizando-o pela capacidade de aprender automaticamente representações hierárquicas a partir de dados brutos, em contraposição à engenharia manual de atributos que caracteriza os métodos clássicos de aprendizado de máquina -- distinção conceitual diretamente aplicável ao espaço de atributos espectrais agregados (984 canais candidatos) utilizado nesta pesquisa, no qual arquiteturas profundas são comparadas com modelos clássicos operando sobre o mesmo espaço de entrada. Especificamente no domínio de monitoramento de condição de máquinas, \citet{zhao2019dl} revisam de forma sistemática a aplicação de arquiteturas de aprendizado profundo -- incluindo redes recorrentes, convolucionais e autoencoders -- ao diagnóstico de falhas e à predição de vida útil remanescente de equipamentos rotativos, corroborando, em um levantamento mais amplo do que o de \citet{hoang2019} restrito a rolamentos, a tendência de crescente adoção dessas arquiteturas no diagnóstico de falhas mecânicas e reforçando a pertinência de sua aplicação ao problema de diagnóstico (RQ1) desta tese.
Como referência de comparação clássica, esta pesquisa emprega três algoritmos amplamente estabelecidos na literatura de aprendizado de máquina supervisionado. O algoritmo \textit{Random Forest}, proposto por \citet{breiman2001}, constrói um comitê de árvores de decisão treinadas sobre subamostras aleatórias dos dados e dos atributos, combinando suas previsões por votação ou média, e permanece, mais de duas décadas após sua publicação, um dos métodos de referência para problemas de classificação com espaços de atributos de alta dimensionalidade, como o espaço espectral de 984 canais utilizado nesta tese. As máquinas de vetores de suporte (\textit{Support Vector Machines} -- SVM), formuladas por \citet{cortes1995}, buscam o hiperplano de máxima margem separando classes no espaço de atributos -- eventualmente transformado por funções de núcleo (\textit{kernel}) para espaços de maior dimensionalidade --, sendo historicamente um dos métodos de maior desempenho em problemas de classificação com número limitado de amostras rotuladas, cenário relevante para a base de dados experimental desta pesquisa. Por fim, o algoritmo XGBoost, proposto por \citet{chen2016xgboost}, implementa o princípio de \textit{gradient boosting} -- construção sequencial de árvores de decisão, cada uma corrigindo os erros residuais do comitê anterior -- de forma escalável e regularizada, e tornou-se um dos algoritmos de maior desempenho empírico em competições e aplicações industriais de aprendizado supervisionado sobre dados tabulares, categoria à qual pertence o espaço de atributos espectrais agregados desta tese.
Do lado das arquiteturas sequenciais profundas comparadas nesta pesquisa para o diagnóstico de falha (RQ1), duas contribuições fundacionais merecem destaque. A rede de memória de curto prazo longa (\textit{Long Short-Term Memory} -- LSTM), proposta por \citet{hochreiter1997}, introduz um mecanismo de portas (\textit{gates}) que regula o fluxo de informação ao longo do tempo, mitigando o problema do desvanecimento do gradiente que limitava as redes neurais recorrentes simples (\textit{SimpleRNN}) em sequências longas, e constitui a arquitetura-base tanto do modelo LSTM quanto do modelo híbrido CNN+LSTM comparados nesta tese, além de fundamentar sua variante bidirecional (BiLSTM), que processa a sequência em ambos os sentidos temporais. Posteriormente, \citet{cho2014} propõem a unidade recorrente com portas (\textit{Gated Recurrent Unit} -- GRU), uma simplificação estrutural da LSTM que reduz o número de portas e, consequentemente, de parâmetros treináveis, mantendo capacidade comparável de captura de dependências temporais -- arquitetura também incluída no conjunto de seis modelos sequenciais comparados nesta pesquisa. Complementarmente, \citet{vaswani2017} introduzem a arquitetura Transformer, que substitui a recorrência por um mecanismo de atenção capaz de ponderar diretamente a relevância entre todos os pontos de uma sequência, independentemente de sua distância temporal, e cuja adaptação a sinais espectrais temporais constitui o sexto modelo sequencial comparado nesta tese, permitindo avaliar se a captura de dependências de longo alcance por atenção supera, para o problema de diagnóstico de falha aqui tratado, as arquiteturas recorrentes com portas.
Por fim, dado que esta pesquisa lida com um espaço de atributos espectrais de alta dimensionalidade (984 canais candidatos) cuja interpretabilidade é relevante tanto para a validação física dos modelos quanto para a comunicação dos resultados a um público de engenharia, esta tese emprega o método SHAP (\textit{SHapley Additive exPlanations}), proposto por \citet{lundberg2017}, que unifica, sob uma base teórica comum fundamentada na teoria de jogos cooperativos (valores de Shapley), diversos métodos de interpretabilidade previamente propostos de forma independente, atribuindo a cada atributo de entrada uma contribuição marginal consistente para a saída do modelo. A aplicação do SHAP ao espaço de atributos espectrais desta pesquisa permite identificar quais ordens harmônicas e interharmônicas específicas mais contribuem para o diagnóstico de falha e para a desagregação de cargas, conectando diretamente os resultados obtidos (Capítulo~\ref{cap:resultados}) à motivação, discutida na Seção~\ref{sec:sintese}, quanto à relevância diagnóstica potencial das componentes interharmônicas evidenciada por \citet{efe2016}.
\subsection{Síntese crítica e lacunas identificadas}
\label{sec:sintese}
A leitura conjunta dos vinte e um trabalhos revisados neste capítulo -- distribuídos entre a literatura de NILM baseado em conteúdo harmônico, a literatura de diagnóstico de falhas via MCSA e os fundamentos dos métodos de aprendizado de máquina e aprendizado profundo empregados na metodologia (Seção~\ref{sec:ml_dl}) -- permite identificar convergências e divergências relevantes para o posicionamento desta tese. Há convergência, atravessando as duas vertentes de literatura diretamente relacionadas ao problema de pesquisa -- NILM baseado em corrente harmônica \citep{hart1992,zeifman2011,djordjevic2017,bouhouras2017,li2019moea} e diagnóstico de falha de máquinas rotativas via MCSA \citep{nandi2005,thomson2001,niu2023,benamira2022,martinez2022} --, quanto à premissa de que o conteúdo espectral da corrente (e, em menor medida, da tensão) já disponível em pontos de medição existentes constitui uma fonte de informação eficaz e computacionalmente mais leve do que a captura e o processamento contínuo da forma de onda bruta ou a instalação de sensoriamento dedicado, seja para identificar cargas elétricas, seja para diagnosticar a condição de uma máquina. Há divergência quanto à extensão do espaço de atributos espectrais considerado: enquanto \citet{djordjevic2017}, \citet{bouhouras2017} e a MCSA clássica de \citet{thomson2001} optam deliberadamente por reduzir esse espaço a um número mínimo de ordens harmônicas ou a bandas laterais específicas, \citet{li2019moea} demonstram, do lado do NILM, que a ampliação do número de atributos tende a melhorar o desempenho de desagregação -- tendência já antecipada, em nível conceitual, por \citet{zeifman2011} -- que esta tese testa, do lado do diagnóstico e da desagregação conjuntos, ao adotar um espaço de 984 canais espectrais candidatos.
Mais relevante ainda para o posicionamento desta tese é uma lacuna que atravessa as duas vertentes de literatura revisadas: nenhum dos doze trabalhos de NILM e de MCSA revisados nas seções anteriores avalia sistematicamente o conteúdo interharmônico como atributo de identificação de carga ou de diagnóstico de falha. A literatura de NILM permanece restrita a ordens harmônicas inteiras \citep{hart1992,zeifman2011,djordjevic2017,bouhouras2017,li2019moea}, apesar de \citet{efe2016} demonstrar, em um sistema retificador estruturalmente análogo ao desta pesquisa, que interharmônicas podem apresentar distorção superior à das harmônicas adjacentes. A literatura clássica de MCSA, por sua vez, concentra-se em bandas laterais associadas a ordens harmônicas inteiras da frequência de escorregamento \citep{thomson2001,nandi2005} ou em indicadores estatísticos do regime transitório \citep{martinez2022}, sem tratar, em nenhum dos casos, de componentes interharmônicas como possível assinatura de falha mecânica. Tampouco há, entre os trabalhos revisados, um sistema que trate diagnóstico de falha mecânica e desagregação de cargas elétricas como duas saídas de um mesmo espaço de entrada espectral agregado em um sistema trifásico industrial -- \citet{benamira2022} e \citet{martinez2022} tratam de falha mecânica isolada ou combinada em cenários específicos, sem desagregação de cargas nem avaliação de conteúdo interharmônico, e os trabalhos de NILM revisados tratam exclusivamente de cargas residenciais monofásicas, sem qualquer componente de diagnóstico de falha mecânica. É precisamente essa dupla lacuna -- a ausência de avaliação sistemática de componentes interharmônicas e a ausência de um sistema unificado de diagnóstico e desagregação em ambiente trifásico industrial -- que motiva as perguntas de pesquisa formuladas na Seção~\ref{sec:objetivos}.

\subsection{Objetivos e perguntas de pesquisa}
\label{sec:objetivos}
A partir da lacuna identificada -- a ausência, na literatura de NILM baseada em corrente harmônica, de uma avaliação sistemática do papel de componentes interharmônicas, bem como a ausência de trabalhos que integrem diagnóstico de falha mecânica e desagregação de cargas elétricas a partir de uma mesma assinatura espectral agregada em sistemas trifásicos --, esta pesquisa formula duas perguntas de pesquisa co-primárias.
\textbf{RQ1 (diagnóstico e manutenção preditiva)}: um modelo treinado sobre a assinatura espectral agregada (harmônicas e interharmônicas de tensão e corrente, sem forma de onda bruta) consegue diagnosticar corretamente o modo de falha de um sistema trifásico -- saudável, falha elétrica isolada, falha mecânica isolada ou falha combinada -- com desempenho suficiente para suportar detecção precoce de falha?
\textbf{RQ2 (desagregação de cargas)}: com duas cargas elétricas operando simultaneamente no mesmo ponto de acoplamento comum, é possível desagregar suas contribuições espectrais individuais utilizando exclusivamente a assinatura espectral agregada, sem recorrer à forma de onda bruta e, em particular, as componentes interharmônicas contribuem de forma relevante para essa desagregação, tal como sugerido pelos resultados de \citet{efe2016}?
%===============================================================
\section{Metodologia}
\label{cap:metodologia}
Este capítulo descreve a metodologia experimental empregada para investigar as duas perguntas de pesquisa formuladas na Seção~\ref{sec:objetivos} do Capítulo~\ref{cap:fundamentacao}: o diagnóstico de falha em máquinas rotativas (RQ1) e a desagregação de cargas elétricas (RQ2), ambas a partir de uma mesma assinatura espectral agregada de tensão e corrente. A Seção~\ref{sec:bancadas} apresenta as bancadas experimentais utilizadas para a coleta de dados. A Seção~\ref{sec:bd} descreve a base de dados resultante e a evolução do escopo de análise adotado ao longo da pesquisa. A Seção~\ref{sec:atributos} detalha o espaço de atributos espectrais extraído dos sinais brutos. A Seção~\ref{sec:arquiteturas} apresenta os modelos de aprendizado de máquina e aprendizado profundo comparados -- em particular, as arquiteturas sequenciais (RNNs e variantes) empregadas no diagnóstico de falha --, cujos fundamentos conceituais já foram revisados na Seção~\ref{sec:ml_dl} do Capítulo~\ref{cap:fundamentacao}. A Seção~\ref{sec:protocolo} descreve o protocolo de treinamento e avaliação, e a Seção~\ref{sec:interpretabilidade} descreve a análise de interpretabilidade via SHAP. Os resultados obtidos a partir desta metodologia são apresentados e discutidos no Capítulo~\ref{cap:resultados}.
\subsection{Bancadas experimentais}
\label{sec:bancadas}
A coleta de dados que fundamenta esta pesquisa foi realizada em duas bancadas experimentais complementares, refletindo a natureza dupla dos objetivos de pesquisa: uma bancada voltada ao diagnóstico de condição de máquina rotativa (RQ1) e uma bancada voltada à geração de cenários de coexistência de cargas elétricas não lineares (RQ2). Ambas as bancadas compartilham o mesmo princípio de instrumentação -- medição trifásica de tensão e corrente em um ponto de acoplamento comum, sem sensoriamento adicional dedicado à máquina ou à carga individual --, de modo a preservar, também em ambiente experimental controlado, a premissa de não intrusividade discutida na Seção~\ref{sec:sintese} do Capítulo~\ref{cap:fundamentacao}. \textcolor{blue}{O instrumento de medição, comum às duas bancadas, é um analisador de qualidade de energia IMS PowerNET PQ-700 G4 (modelo PQK7-60-6C7J-D1N-0-3; IMS -- Soluções em Energia Ltda., Porto Alegre/RS, Brasil -- Figura~\ref{fig:placa_analisador}), com comunicação USB, 512~MB de memória interna, sistema trifásico Y/$\Delta$, frequência de rede de 60~Hz, consumo de 20~VA, faixa de tensão de 50--300~V\textsubscript{ca} (fase-neutro) / 50--500~V\textsubscript{ca} (fase-fase) e sonda de corrente flexível Flex~1000-e. Os dados são extraídos do analisador pelo software proprietário PowerMANAGER (IMS). \textcolor{blue}{Segundo o manual técnico da série PowerNET PQ-700 (IMS Soluções em Energia), a medição de harmônicos e inter-harmônicos segue a norma IEC~61000-4-7, com janela de análise de 200~ms, resolução espectral de 5~Hz, ordens harmônicas e inter-harmônicas de 0 a 50, e agrupamento configurável entre Grupo, Subgrupo e Nenhum; distúrbios de tensão são medidos conforme IEC~61000-4-30, e \textit{flicker} conforme IEC~61000-4-15. A aquisição é simultânea em todos os canais, com resolução de 16~bits e taxa de amostragem de 10.240 amostras por segundo. O agrupamento efetivamente configurado neste equipamento, verificado diretamente na tela do analisador, é Subgrupo (S-Grupo).}}

\begin{figure}[H]
    \centering
    \fbox{%
        \begin{minipage}[c][6cm][c]{0.7\textwidth}
            \centering
            \textit{\textcolor{blue}{Espaço reservado para a foto da placa de identificação do analisador IMS PowerNET PQ-700~G4 (salvar em \texttt{figuras/placa\_analisador.jpg})}}
        \end{minipage}%
    }
    \caption{\textcolor{blue}{Placa de identificação do analisador de qualidade de energia IMS PowerNET PQ-700~G4, modelo PQK7-60-6C7J-D1N-0-3.}}
    \label{fig:placa_analisador}
\end{figure}
\subsubsection{Bancada de diagnóstico de falhas (RQ1)}
\label{subsec:bancada_rq1}
A primeira bancada reproduz um motor de indução trifásico instrumentado, operado sob as condições saudável, falha elétrica isolada, falha mecânica isolada (desbalanceamento e desalinhamento) e falha combinada (elétrica + mecânica), conforme a taxonomia experimental adotada nesta tese e já introduzida no Capítulo~\ref{cap:fundamentacao}. \textcolor{blue}{O motor de indução trifásico utilizado é um WEG W22 Plus (Figura~\ref{fig:placa_motor}), potência nominal de 0,75~kW (1,0~cv), rotação nominal de 1.730~rpm (4 polos, 60~Hz), tensão dupla 220/380~V (ligação $\Delta$/Y, respectivamente), corrente nominal de 3,02~A (220~V, $\Delta$) / 1,75~A (380~V, Y), grau de proteção IP55, regime de serviço contínuo (S1), rolamentos 6204-ZZ (dianteiro) e 6203-ZZ.EM (traseiro) -- dados extraídos diretamente da placa de identificação do motor. \textcolor{blue}{O campo de massa (Kg) da placa está em branco (não preenchido pelo fabricante); se a massa for necessária para o apêndice de reprodutibilidade, deve ser obtida por pesagem direta do motor ou consultada no catálogo WEG para o modelo W22 Plus 0,75~kW/1,0~cv. O acionamento é por partida direta. Como a corrente nominal do motor (3,02~A em 220~V) é muito pequena frente à faixa do sensor de corrente flexível Flex~1000-e (dimensionado para até 1000~A), o condutor foi enrolado três vezes através do sensor para multiplicar por três a corrente efetivamente percebida pela bobina de Rogowski, aumentando a resolução da medição. Essa multiplicação não é compensada na parametrização do analisador nem em pós-processamento: os canais de corrente desta bancada são utilizados em valor bruto, na escala de saída do sensor (aproximadamente três vezes a corrente real no condutor), não na escala de corrente física em Ampères do motor. Essa escolha não compromete a validade dos resultados de classificação e sensibilidade desta pesquisa, que dependem de padrões relativos entre classes dentro de um mesmo canal, não de calibração absoluta; deve, porém, ser observada ao comparar qualquer valor de corrente reportado nesta bancada com os dados de placa do motor (Ampères físicos). [Pendente de confirmação: se a mesma técnica de multiplicação de espiras foi empregada na instrumentação de corrente da bancada de retificador (RQ2, Seção~\ref{subsec:bancada_rq2}) -- caso os fatores de multiplicação sejam diferentes entre as duas bancadas, isso deve ser declarado explicitamente, já que RQ2 combina corrente elétrica de fonte com corrente da máquina no mesmo canal.]}}
\begin{figure}[htbp]
    \centering
    \includegraphics[width=0.8\textwidth]{figuras/Bancada Motomáquina 1.jpeg}
    \fbox{%
    }
    \caption{Bancada experimental utilizada para a coleta de dados de diagnóstico de falha de máquina rotativa (RQ1), instrumentada para aquisição trifásica de tensão e corrente no ponto de acoplamento comum.}
    \label{fig:bancada_rq1}
\end{figure}
\begin{figure}[H]
    \centering
    \fbox{%
        \begin{minipage}[c][6cm][c]{0.7\textwidth}
            \centering
            \textit{\textcolor{blue}{Espaço reservado para a foto da placa de identificação do motor WEG W22 Plus (salvar em \texttt{figuras/placa\_motor.jpg})}}
        \end{minipage}%
    }
    \caption{\textcolor{blue}{Placa de identificação do motor de indução trifásico WEG W22 Plus utilizado na bancada de diagnóstico de falhas (RQ1).}}
    \label{fig:placa_motor}
\end{figure}
\subsubsection{Bancada de retificador trifásico (RQ2)}
\label{subsec:bancada_rq2}
A segunda bancada reproduz um retificador trifásico alimentando cargas não lineares controladas, permitindo compor combinações de cargas operando simultaneamente no mesmo ponto de acoplamento comum -- topologia estruturalmente análoga à utilizada por \citet{efe2016} para caracterizar o comportamento interharmônico de sistemas retificadores (Seção~\ref{sec:sintese}, Capítulo~\ref{cap:fundamentacao}) e que fundamenta a hipótese, testada nesta tese, de que componentes interharmônicas carregam informação diagnóstica relevante também para a desagregação de cargas (RQ2). [Inserir aqui: topologia exata do retificador, cargas não lineares utilizadas e critério de composição das combinações de carga.]
\begin{figure}[htbp]
    \centering
    \includegraphics[width=0.8\textwidth]{figuras/Retificador Trifásico.jpeg}
    \fbox{%
    }
    \caption{Bancada experimental utilizada para a coleta de dados de desagregação de cargas (RQ2), baseada em retificador trifásico alimentando combinações controladas de cargas não lineares.}
    \label{fig:bancada_rq2}
\end{figure}
\subsubsection{Esquemático das bancadas}
\label{subsec:esquematico}
A Figura~\ref{fig:esquematico} sintetiza a topologia de instrumentação comum às duas bancadas descritas nas Seções~\ref{subsec:bancada_rq1} e~\ref{subsec:bancada_rq2}: aquisição trifásica de tensão e corrente no ponto de acoplamento comum, seguida de processamento espectral (decomposição em ordens harmônicas e interharmônicas, conforme detalhado na Seção~\ref{sec:atributos}) e alimentação do espaço de atributos resultante aos modelos de classificação e desagregação descritos na Seção~\ref{sec:arquiteturas}.
\begin{figure}[htbp]
    \centering
    % \includegraphics[width=0.9\textwidth]{figuras/esquematico_bancadas.png}
    \fbox{%
        \begin{minipage}[c][7cm][c]{0.9\textwidth}
            \centering
            \textit{Espaço reservado para esquemático das bancadas experimentais\\
            (pontos de medição, sensores, cadeia de aquisição e processamento espectral)}
        \end{minipage}%
    }
    \caption{Esquemático da instrumentação das bancadas experimentais, com indicação do ponto de acoplamento comum, dos sensores de tensão e corrente trifásicos e da cadeia de processamento espectral que alimenta o espaço de atributos descrito na Seção~\ref{sec:atributos}.}
    \label{fig:esquematico}
\end{figure}
\subsection{Base de dados}
\label{sec:bd}
A base de dados desta pesquisa foi construída a partir das bancadas descritas na Seção~\ref{sec:bancadas} e passou por uma reformulação metodológica relevante ao longo do desenvolvimento da tese, que é registrada nesta seção por sua relevância direta para a interpretação dos resultados apresentados no Capítulo~\ref{cap:resultados}.
Na concepção inicial do projeto, o escopo de análise foi deliberadamente restrito às componentes de corrente harmônica e interharmônica, seguindo diretamente a linha de investigação de \citet{djordjevic2017} e \citet{bouhouras2017} (Seção~\ref{sec:sintese}, Capítulo~\ref{cap:fundamentacao}) e a motivação específica de \citet{efe2016} quanto à relevância diagnóstica potencial das interharmônicas. Essa delimitação inicial cumpriu um papel metodológico importante: permitiu validar, em um espaço de atributos reduzido e diretamente interpretável, a hipótese de que o conteúdo espectral de corrente carrega informação diagnosticamente relevante, tanto para falha de máquina (RQ1) quanto para composição de carga (RQ2), antes de se investir na ampliação do escopo de análise.
Em uma etapa subsequente, o escopo de análise foi ampliado para incluir também as componentes de tensão -- e não apenas de corrente -- e para deixar de restringir, \textit{a priori}, o conjunto de atributos considerado relevante. Em vez de selecionar manualmente um subconjunto de ordens harmônicas e interharmônicas com base em hipóteses formuladas previamente, optou-se por disponibilizar aos modelos de aprendizado de máquina e aprendizado profundo (Seção~\ref{sec:arquiteturas}) o espaço espectral candidato de forma ampla, permitindo que os próprios modelos -- e, posteriormente, a análise de interpretabilidade via SHAP (Seção~\ref{sec:interpretabilidade}) -- identificassem quais canais espectrais são efetivamente relevantes para cada tarefa, em vez de impor essa seleção antecipadamente. Essa mudança de escopo é diretamente respaldada pela literatura revisada no Capítulo~\ref{cap:fundamentacao}: \citet{li2019moea} demonstram, do lado do NILM, que a ampliação do número de atributos considerados tende a melhorar o desempenho de desagregação, achado também antecipado, em nível conceitual, por \citet{zeifman2011}; a presente reformulação metodológica pode, portanto, ser lida como um teste direto dessa tendência da literatura, estendida simultaneamente ao problema de diagnóstico de falha (RQ1).
\textcolor{blue}{O intervalo de registro (\textit{datalog}) configurado no analisador é de 10~s entre aquisições consecutivas -- cada aquisição sendo, por sua vez, uma agregação interna do equipamento sobre uma janela de 200~ms (Seção~\ref{sec:atributos}). É esse intervalo de 10~s entre registros, não a janela de 200~ms, que determina o número de aquisições por sessão experimental (Seção~\ref{sec:bancadas}).} [Inserir aqui: número total de amostras, distribuição entre classes/condições de falha e entre combinações de carga, e eventual divisão em subconjuntos de treino, validação e teste, esta última também detalhada na Seção~\ref{sec:protocolo}.]
\subsection{Espaço de atributos espectrais}
\label{sec:atributos}
O espaço de atributos resultante da etapa de processamento descrita na Seção~\ref{sec:bd} compreende 984 canais espectrais candidatos, obtidos a partir da decomposição harmônica e interharmônica das três fases de tensão e corrente medidas no ponto de acoplamento comum de cada bancada (Seção~\ref{sec:bancadas}). Esse espaço de atributos constitui a entrada comum a todos os modelos comparados na Seção~\ref{sec:arquiteturas}, tanto para o diagnóstico de falha (RQ1) quanto para a desagregação de cargas (RQ2), permitindo uma comparação direta de desempenho entre arquiteturas sobre um mesmo conjunto de dados de entrada. \textcolor{blue}{Segundo o manual técnico da série PowerNET PQ-700 (IMS Soluções em Energia), a extração espectral segue a norma IEC~61000-4-7, com janela de análise de 200~ms, aquisição simultânea de todos os canais a 16~bits e taxa de amostragem de 10.240 amostras por segundo, resolução espectral de 5~Hz e ordens harmônicas e inter-harmônicas de 0 a 50. O agrupamento espectral utilizado, confirmado diretamente na tela do equipamento, é Subgrupo (S-Grupo) -- ou seja, cada componente central (harmônico) é agrupado ao componente espectral imediatamente anterior e posterior. O método de extração (FFT interna ao analisador) é proprietário do equipamento, não implementado pelo pipeline desta pesquisa; os 984 canais candidatos citados acima são a saída já processada pelo analisador, exportada via PowerMANAGER.}

\textcolor{blue}{\subsubsection{Definição metrológica das componentes harmônicas e inter-harmônicas}
\label{sec:definicao_metrologica}
Para que um nome de canal (por exemplo, $H_{05}\_I_b\_avg$ ou $IH_{46}\_I_c\_avg$) seja interpretável fisicamente, e não apenas como identificador categórico de coluna, esta seção formaliza o mapeamento canal~$\to$~frequência~$\to$~grandeza~$\to$~unidade subjacente aos 984 canais candidatos, respondendo diretamente à observação do orientador de que essa definição não pode ser inferida apenas do nome da planilha de origem. Sob a norma IEC~61000-4-7 já citada, com frequência fundamental $f_1=60$~Hz: a componente harmônica de ordem $n$ corresponde à frequência $n \cdot f_1$ (por exemplo, $H_{05}$ corresponde a $300$~Hz); a componente inter-harmônica de ordem $n$, sob o método de subgrupo (S-Grupo) confirmado na tela do equipamento (Seção~\ref{sec:bancadas}), corresponde ao subgrupo espectral entre as ordens harmônicas $n$ e $n+1$, ou seja, à banda $(n \cdot f_1,\ (n{+}1) \cdot f_1)$ Hz (por exemplo, $IH_{46}$ corresponde à banda entre $2.760$~Hz e $2.820$~Hz) -- esta é a definição padrão do subgrupo inter-harmônico na norma, não uma convenção específica desta pesquisa. A grandeza de cada canal é a amplitude espectral (tensão ou corrente, conforme o sufixo $V$/$I$ e a fase $a$/$b$/$c$ no nome do canal), expressa como percentual da componente fundamental medida na mesma fase -- unidade já implícita nos valores percentuais reportados ao longo deste capítulo (por exemplo, Seção~\ref{sec:resultados_bd}). \textbf{Dois pontos, porém, permanecem pendentes de confirmação direta} e não devem ser lidos como estabelecidos por este mapeamento: primeiro, se a indexação de ordem inter-harmônica realmente utilizada pelo firmware do PowerNET/exportação do PowerMANAGER segue estritamente essa convenção de subgrupo da norma, ou uma variante proprietária do fabricante; segundo, se os sufixos \texttt{\_avg} e \texttt{\_max} correspondem a agregações já calculadas internamente pelo equipamento dentro de cada registro de 10~s (Seção~\ref{sec:bd}), ou a estatísticas calculadas posteriormente pelo pipeline desta pesquisa sobre a série de aquisições de uma sessão -- distinção que, embora não afete a validade dos resultados de classificação e regressão desta tese (que dependem de padrões relativos entre canais e classes, não de calibração absoluta em unidades físicas), deve ser esclarecida junto ao fabricante ou à documentação completa do PowerMANAGER antes de qualquer leitura em unidades físicas absolutas (Ampères, Volts) dos valores percentuais reportados.}

\subsection{Modelos comparados}
\label{sec:arquiteturas}
Esta seção apresenta os modelos de aprendizado de máquina e aprendizado profundo comparados nesta pesquisa, cujos fundamentos conceituais e históricos foram revisados na Seção~\ref{sec:ml_dl} do Capítulo~\ref{cap:fundamentacao}. Os modelos são apresentados nesta seção -- antes da exposição dos resultados obtidos, no Capítulo~\ref{cap:resultados} -- de modo que o leitor disponha do panorama completo das arquiteturas comparadas ao interpretar o desempenho relativo entre elas.
\subsubsection{Modelos clássicos de referência}
\label{subsec:modelos_classicos}
\textcolor{blue}{Seis modelos são empregados como referência de comparação tabular para os modelos de aprendizado profundo sequencial descritos na Seção~\ref{subsec:rnns}: \textit{Random Forest} \citep{breiman2001}, Máquinas de Vetores de Suporte (SVM) \citep{cortes1995} e XGBoost \citep{chen2016xgboost} -- já descritos na Seção~\ref{sec:ml_dl} --, acrescidos de um Perceptron Multicamadas (MLP, duas camadas ocultas), do algoritmo de $k$ vizinhos mais próximos ($k$-NN, $k=5$) e de Regressão Logística. Os seis modelos operam diretamente sobre o espaço de atributos espectrais descrito na Seção~\ref{sec:atributos}, tratado como vetor de atributos estático, sem exploração explícita de eventual estrutura sequencial entre canais espectrais. Cabe notar que o MLP, embora seja tecnicamente uma rede neural, não possui mecanismo de recorrência nem de atenção; por isso é tratado aqui como referência tabular, e não como parte da comparação de arquiteturas sequenciais da Seção~\ref{subsec:rnns}.}
\subsubsection{Arquiteturas sequenciais de aprendizado profundo (RNNs e variantes)}
\label{subsec:rnns}
Para o problema de diagnóstico de falha (RQ1), seis arquiteturas sequenciais de aprendizado profundo foram implementadas e comparadas sistematicamente, endereçando diretamente a lacuna apontada por \citet{hoang2019} quanto à escassez de redes neurais recorrentes aplicadas ao diagnóstico de falha de máquinas rotativas (Seção~\ref{sec:sintese}, Capítulo~\ref{cap:fundamentacao}). A Tabela~\ref{tab:arquiteturas} sintetiza as seis arquiteturas, seu mecanismo central e a referência fundacional correspondente, já discutida em detalhe na Seção~\ref{sec:ml_dl} do Capítulo~\ref{cap:fundamentacao}.
\begin{table}[htbp]
    \centering
    \caption{Arquiteturas sequenciais de aprendizado profundo comparadas para o diagnóstico de falha (RQ1).}
    \label{tab:arquiteturas}
    \begin{tabular}{@{}p{2.6cm}p{5.3cm}p{2.2cm}p{3cm}@{}}
        \toprule
        \textbf{Arquitetura} & \textbf{Mecanismo central} & \textbf{Referência} & \textbf{Papel na comparação} \\
        \midrule
        SimpleRNN & Recorrência simples, sem mecanismo de portas & --- & Linha de base recorrente \\
        LSTM & Portas de entrada/esquecimento/saída, mitigação do desvanecimento do gradiente & \citet{hochreiter1997} & Arquitetura recorrente de referência \\
        GRU & Simplificação da LSTM, portas de atualização/redefinição & \citet{cho2014} & Alternativa recorrente mais leve \\
        BiLSTM & LSTM bidirecional, processamento em ambos os sentidos temporais & \citet{hochreiter1997} & Extensão bidirecional da LSTM \\
        CNN+LSTM & Extração convolucional de padrões locais seguida de recorrência LSTM & \citet{hochreiter1997} & Modelo híbrido convolucional-recorrente \\
        Transformer & Mecanismo de atenção, sem recorrência & \citet{vaswani2017} & Captura de dependências de longo alcance \\
        \bottomrule
    \end{tabular}
\end{table}
Cabe registrar que a arquitetura SimpleRNN, apesar de não possuir uma referência fundacional isolada citável -- por corresponder à formulação mais elementar de rede neural recorrente, historicamente anterior aos mecanismos de porta introduzidos por \citet{hochreiter1997} --, é incluída na comparação como linha de base recorrente, permitindo isolar, nos resultados apresentados no Capítulo~\ref{cap:resultados}, a contribuição específica de cada mecanismo adicional (portas, bidirecionalidade, convolução, atenção) em relação à recorrência pura.
A Figura~\ref{fig:fluxograma_rnns} sintetiza o fluxo de modelagem comparativa descrito nesta seção, desde o espaço de atributos espectrais de entrada (Seção~\ref{sec:atributos}) até a avaliação comparativa das seis arquiteturas sequenciais e a análise de interpretabilidade via SHAP (Seção~\ref{sec:interpretabilidade}), cujos resultados são apresentados no Capítulo~\ref{cap:resultados}. O fluxograma evidencia que as seis arquiteturas são treinadas e avaliadas de forma paralela e independente sobre o mesmo espaço de atributos de entrada e sob o mesmo protocolo de divisão de dados (Seção~\ref{sec:protocolo}), condição necessária para que as diferenças de desempenho observadas nos resultados possam ser atribuídas às arquiteturas em si, e não a diferenças no pré-processamento ou no particionamento dos dados.
\begin{figure}[htbp]
\centering
\begin{tikzpicture}[
    node distance=0.9cm and 0.55cm,
    block/.style={draw, rounded corners, fill=gray!10, align=center, minimum height=1cm, text width=5.6cm, font=\small},
    model/.style={draw, rounded corners, fill=blue!8, align=center, minimum height=0.9cm, text width=2.1cm, font=\footnotesize},
    arrow/.style={-{Latex[length=2mm]}, thick}
]
\node[block] (input) {Espaço de atributos espectrais (984 canais) -- Seção~\ref{sec:atributos}};
\node[block, below=1.0cm of input] (prep) {Pré-processamento e divisão treino/validação/teste -- Seção~\ref{sec:protocolo}};
\node[model, below=1.6cm of prep, xshift=-6.3cm] (m1) {SimpleRNN};
\node[model, right=of m1] (m2) {LSTM};
\node[model, right=of m2] (m3) {GRU};
\node[model, right=of m3] (m4) {BiLSTM};
\node[model, right=of m4] (m5) {CNN+LSTM};
\node[model, right=of m5] (m6) {Transformer};
\coordinate (busTop) at ($(prep.south)+(0,-0.6)$);
\draw[thick] (prep.south) -- (busTop);
\draw[thick] (m1.north |- busTop) -- (m6.north |- busTop);
\foreach \m in {m1,m2,m3,m4,m5,m6}{
  \draw[arrow] (\m.north |- busTop) -- (\m.north);
}
\coordinate (busBotMid) at ($(m1.south)!0.5!(m6.south)+(0,-0.7)$);
\draw[thick] (m1.south |- busBotMid) -- (m6.south |- busBotMid);
\foreach \m in {m1,m2,m3,m4,m5,m6}{
  \draw[thick] (\m.south) -- (\m.south |- busBotMid);
}
\node[block, below=1.1cm of busBotMid] (eval) {Avaliação comparativa (métricas de classificação) -- Seção~\ref{sec:protocolo}};
\draw[arrow] (busBotMid) -- (eval.north);
\node[block, below=of eval] (shap) {Interpretabilidade via SHAP -- Seção~\ref{sec:interpretabilidade}};
\draw[arrow] (eval) -- (shap);
\node[block, below=of shap] (result) {Resultados de diagnóstico de falha (RQ1) -- Capítulo~\ref{cap:resultados}};
\draw[arrow] (shap) -- (result);
\end{tikzpicture}
\caption{Fluxograma da modelagem comparativa das seis arquiteturas sequenciais de aprendizado profundo empregadas no diagnóstico de falha (RQ1), do espaço de atributos espectrais de entrada até a avaliação comparativa e a análise de interpretabilidade.}
\label{fig:fluxograma_rnns}
\end{figure}
\subsection{Protocolo de treinamento e avaliação}
\label{sec:protocolo}
\textcolor{blue}{Cada uma das 17 condições experimentais corresponde a uma única sessão de aquisição contínua (Seção~\ref{sec:bd}), não a múltiplas execuções independentes -- não há, portanto, uma coluna de grupo que tornasse \texttt{GroupKFold} informativo (grupo coincidiria com rótulo). Um particionamento aleatório estratificado, como usado em uma versão preliminar deste protocolo, testa generalização entre aquisições vizinhas da mesma sessão, não entre execuções experimentais independentes, e superestima o desempenho por isso (quantificado na Seção~\ref{sec:resultados_rq1}). O protocolo adotado aqui particiona cada rótulo cronologicamente: as aquisições de cada sessão são ordenadas por \textit{timestamp} e a fração final da série temporal de cada rótulo é reservada para teste (\textit{holdout} único, 70/30) ou para uma das 5 dobras (validação cruzada), com as demais aquisições daquele rótulo, sempre anteriores no tempo, formando o treino/dobras restantes -- um K-fold temporal em blocos contíguos por rótulo, análogo ao \textit{holdout} cronológico, em vez do \texttt{StratifiedKFold} aleatório usado antes. A seleção de atributos (Seção~\ref{subsec:modelos_classicos}) é reajustada a cada \textit{split}/dobra usando exclusivamente as linhas de treino dessa dobra -- o F-score ANOVA e a eliminação de redundância por correlação não veem, em nenhuma etapa, as linhas reservadas para teste daquela dobra. Essa correção também se aplica a RQ2 (Seção~\ref{sec:resultados_rq2}): ali a seleção de atributos é reajustada excluindo as quatro condições combinadas inteiras -- não apenas a condição retida em cada dobra do \textit{leave-one-condition-out} --, de modo que a mesma lista de 60 canais seja usada em todas as dobras e nenhuma condição combinada influencie, em momento algum, quais canais são selecionados. O escalonamento (\texttt{StandardScaler}, para os modelos sensíveis a escala) já ocorre dentro do \textit{pipeline} de cada modelo (Seção~\ref{subsec:modelos_classicos}), portanto também é reajustado por \textit{split}. Métricas: acurácia e F1-macro (que pondera igualmente todas as classes, relevante dado o desbalanceamento entre condições -- Seção~\ref{sec:resultados_bd}) para os seis modelos clássicos, complementadas por matriz de confusão. Hiperparâmetros: valores padrão de cada implementação \textit{scikit-learn}/\textit{XGBoost} (Seção~\ref{subsec:modelos_classicos}), sem busca de hiperparâmetros na comparação principal -- a busca aleatória (Seção~\ref{subsec:rnns}) é demonstrada separadamente, apenas para a LSTM. As seis arquiteturas sequenciais (Seção~\ref{subsec:rnns}) foram retreinadas sob o mesmo particionamento temporal: a divisão 75/25 por rótulo é aplicada às aquisições \emph{antes} do janelamento, e as janelas deslizantes (\textit{stride}=4 sobre janelas de tamanho~8) são geradas separadamente dentro de cada metade -- não sobre o conjunto completo seguido de divisão aleatória das janelas, como na versão preliminar deste protocolo --, de modo que nenhuma janela de treino compartilha aquisições com uma janela de validação, mesmo havendo sobreposição interna a cada metade (verificado por amostragem: Seção~\ref{sec:resultados_rnn}). As sementes de treinamento continuam únicas por arquitetura (ver Seção~\ref{cap:limitacoes}) -- essa correção ainda está pendente.}
\subsection{Interpretabilidade via SHAP}
\label{sec:interpretabilidade}
Conforme fundamentado na Seção~\ref{sec:ml_dl} do Capítulo~\ref{cap:fundamentacao}, o método SHAP \citep{lundberg2017} é aplicado ao espaço de atributos espectrais descrito na Seção~\ref{sec:atributos}, com o objetivo de identificar quais ordens harmônicas e interharmônicas, de tensão e de corrente, mais contribuem para as decisões dos modelos comparados na Seção~\ref{sec:arquiteturas}, tanto no diagnóstico de falha (RQ1) quanto na desagregação de cargas (RQ2). Essa análise permite conectar diretamente os resultados obtidos (Capítulo~\ref{cap:resultados}) à motivação teórica quanto à relevância diagnóstica potencial das componentes interharmônicas, evidenciada por \citet{efe2016} e discutida na Seção~\ref{sec:sintese} do Capítulo~\ref{cap:fundamentacao}, avaliando empiricamente se essa relevância se confirma nos modelos treinados sobre os dados das bancadas descritas na Seção~\ref{sec:bancadas}.
%===============================================================
\section{Resultados}
\label{cap:resultados}

%% PROVENIÊNCIA DOS DADOS: todos os números, tabelas e figuras desta seção
%% vêm diretamente do relatório consolidado do autor (relatorio_consolidado.tex),
%% dos arquivos JSON de dados do pipeline (findings.json, tecnico.json, rnn.json,
%% full_report.json) e da inspeção visual de cada uma das 25 figuras em PDF geradas
%% pelo próprio pipeline (harmonic_signatures), lidos diretamente do computador do
%% autor. Nenhum valor foi inventado ou estimado; ressalvas metodológicas já
%% registradas no relatório original (ex.: degenerescência do R^2 por dobra em RQ2,
%% Seção~\ref{sec:resultados_rq2}) foram preservadas e explicadas, não silenciadas.
%% [CORRIGIDO] O parágrafo de referência a "Seção~\ref{sec:taxonomia}" (na discussão
%% da matriz de confusão de 17 classes) apontava para um rótulo nunca definido;
%% corrigido para "Seção~\ref{sec:resultados_bd}", onde a taxonomia de fato aparece
%% (Tabela~\ref{tab:res_taxonomia}).

Este capítulo apresenta os resultados obtidos com a metodologia descrita no Capítulo~\ref{cap:metodologia}, organizados em torno das duas perguntas de pesquisa formalizadas na Seção~\ref{sec:objetivos} do Capítulo~\ref{cap:fundamentacao}. A Seção~\ref{sec:resultados_bd} apresenta a composição final da base de dados e as verificações de qualidade realizadas antes de qualquer modelagem. A Seção~\ref{sec:resultados_rq1} reúne a evidência primária de RQ1 (diagnóstico de falha nas 17 condições experimentais). A Seção~\ref{sec:resultados_sensibilidade} apresenta a sub-análise de sensibilidade mecânica isolada, cujo papel é interpretativo -- não uma terceira pergunta de pesquisa independente --, fortalecendo a leitura das classes mecânicas dentro do diagnóstico geral. A Seção~\ref{sec:resultados_rq2} reúne a evidência primária de RQ2 (desagregação de cargas). A Seção~\ref{sec:resultados_rnn} detalha o treinamento das seis arquiteturas sequenciais introduzidas na Seção~\ref{subsec:rnns} do Capítulo~\ref{cap:metodologia}. A Seção~\ref{sec:resultados_robustez} avalia a robustez dos classificadores por validação cruzada, matrizes de confusão e projeção em componentes principais. A interpretação crítica destes resultados à luz da literatura revisada no Capítulo~\ref{cap:fundamentacao} é aprofundada no Capítulo~\ref{cap:discussao}; este capítulo já introduz, pontualmente, as conexões mais diretas entre achado e literatura, sem antecipar a discussão completa.

\subsection{Base de dados: composição e verificação de qualidade}
\label{sec:resultados_bd}

A base de dados experimental resultante da coleta descrita na Seção~\ref{sec:bancadas} totaliza 19.501 aquisições, distribuídas em 17 condições experimentais mutuamente identificáveis, particionadas em três categorias mutuamente exclusivas para fins de análise (Tabela~\ref{tab:res_taxonomia}). Todas as 17 condições compõem o problema de diagnóstico geral (RQ1); o subconjunto de condições mecânicas puras é reaproveitado à parte, na Seção~\ref{sec:resultados_sensibilidade}, para a sub-análise de sensibilidade isolada.

\begin{table}[H]
    \centering
    \caption{Taxonomia experimental das 17 condições da base de dados.}
    \label{tab:res_taxonomia}
    \begin{tabular}{@{}lrrl@{}}
        \toprule
        \textbf{Categoria} & \textbf{N. condições} & \textbf{N. amostras} & \textbf{Papel na pesquisa} \\
        \midrule
        Mecânica pura & 6 & 8.814 & RQ1 (diagnóstico) + sensibilidade isolada \\
        Elétrica pura & 7 & 4.580 & RQ1 (diagnóstico) + referência elétrica para RQ2 \\
        Combinada (mecânica + elétrica) & 4 & 6.107 & RQ1 (diagnóstico) + RQ2 (desagregação) \\
        \midrule
        \textbf{Total} & \textbf{17} & \textbf{19.501} & --- \\
        \bottomrule
    \end{tabular}
\end{table}

A Figura~\ref{fig:res_labels} mostra a distribuição de amostras entre as 17 condições: as \textcolor{blue}{dez} condições associadas à bancada moto-máquina concentram entre 1.440 e 1.589 amostras cada, enquanto as seis condições associadas ao retificador (\texttt{DIODO~1}--\texttt{6}) e a condição \texttt{RETIFICADOR COMPLETO} somam entre 479 e 736 amostras -- reflexo direto da duração de cada sessão de ensaio, não de um desequilíbrio amostral introduzido pelo pipeline. Essa assimetria é relevante para a leitura dos resultados de classificação da Seção~\ref{sec:resultados_rq1}, e justifica o uso de F1-macro (que pondera igualmente todas as classes) como métrica complementar à acurácia.

\textcolor{blue}{A Tabela~\ref{tab:res_taxonomia_hierarquica} individualiza essas mesmas 17 condições, uma por linha, nos eixos classe, equipamento de origem, estado elétrico, estado mecânico, severidade e categoria científica -- respondendo diretamente à observação do orientador de que a taxonomia experimental, tal como apresentada até então, misturava condição experimental, tipo de falha e severidade dentro do mesmo rótulo plano de 17 classes, sem separar esses três conceitos. A coluna ``Categoria científica'' corresponde à condição de primeiro nível (saudável, falha elétrica isolada, falha mecânica isolada, falha combinada); as colunas ``Estado elétrico'' e ``Estado mecânico'' tornam explícito, para cada uma das 17 classes, que a condição combinada mantém a máquina mecanicamente saudável (balanceada e alinhada) enquanto varia apenas o estado elétrico -- e não uma combinação livre de severidades mecânica e elétrica, como o rótulo plano por si só poderia sugerir. Essa reorganização é adotada como lente de leitura hierárquica dos resultados de RQ1 ao longo do Capítulo~\ref{cap:resultados} e da Discussão (Capítulo~\ref{cap:discussao}): o classificador de primeiro nível é lido na matriz de confusão completa de 17 classes (Figura~\ref{fig:res_cm17}), agregada nos quatro ramos de primeiro nível; o segundo nível já está parcialmente coberto pela sub-análise de sensibilidade mecânica isolada (Seção~\ref{sec:resultados_sensibilidade}), que restringe a leitura ao ramo mecânico e sua escala de severidade (\texttt{RET~N~PARAF}). Cabe uma ressalva importante de escopo: essa reorganização é, nesta versão do trabalho, uma lente de \emph{leitura e relato} dos resultados já obtidos pelo classificador plano de 17 classes -- não a substituição desse classificador por uma arquitetura hierárquica de dois estágios (classificador de ramo seguido de classificador de subtipo por ramo), que permanece como refinamento metodológico de trabalho futuro (Seção~\ref{cap:limitacoes}). Para o ramo elétrico, a semântica exata da numeração \texttt{DIODO~1}--\texttt{6} (se identifica a posição do diodo afetado no retificador trifásico, ou uma contagem ordinal de diodos simultaneamente afetados) depende da topologia exata da bancada de retificador, ainda pendente de detalhamento na Seção~\ref{subsec:bancada_rq2}; por isso a Tabela~\ref{tab:res_taxonomia_hierarquica} lista essas seis condições como identificadores distintos, não como uma escala de severidade confirmada -- diferente do ramo mecânico, no qual a escala \texttt{RET~N~PARAF} já está documentada e discutida (Seção~\ref{sec:discussao_rq1}).}

\begin{table}[H]
    \centering
    \footnotesize
    \caption{\textcolor{blue}{Taxonomia individualizada das 17 condições experimentais, nos eixos condição experimental, equipamento de origem, estado elétrico, estado mecânico, severidade e categoria científica (nível 1 da leitura hierárquica).}}
    \label{tab:res_taxonomia_hierarquica}
    \begin{tabular}{@{}p{3.4cm}p{1.9cm}p{2.0cm}p{2.3cm}p{1.7cm}p{1.7cm}r@{}}
        \toprule
        \textbf{Classe} & \textbf{Equipamento} & \textbf{Estado elétrico} & \textbf{Estado mecânico} & \textbf{Severidade} & \textbf{Categoria científica} & \textbf{N.} \\
        \midrule
        \textcolor{blue}{Máquina balanceada e alinhada} & \textcolor{blue}{Moto-máquina (RQ1)} & \textcolor{blue}{Saudável} & \textcolor{blue}{Saudável} & \textcolor{blue}{---} & \textcolor{blue}{Saudável} & \textcolor{blue}{1.440} \\
        \addlinespace
        \textcolor{blue}{Diodo~1} & \textcolor{blue}{Retificador (RQ2)} & \textcolor{blue}{Falha (diodo~1)} & \textcolor{blue}{Saudável} & \textcolor{blue}{Não confirmada\textsuperscript{a}} & \textcolor{blue}{Elétrica isolada} & \textcolor{blue}{659} \\
        \textcolor{blue}{Diodo~2} & \textcolor{blue}{Retificador (RQ2)} & \textcolor{blue}{Falha (diodo~2)} & \textcolor{blue}{Saudável} & \textcolor{blue}{Não confirmada\textsuperscript{a}} & \textcolor{blue}{Elétrica isolada} & \textcolor{blue}{664} \\
        \textcolor{blue}{Diodo~3} & \textcolor{blue}{Retificador (RQ2)} & \textcolor{blue}{Falha (diodo~3)} & \textcolor{blue}{Saudável} & \textcolor{blue}{Não confirmada\textsuperscript{a}} & \textcolor{blue}{Elétrica isolada} & \textcolor{blue}{648} \\
        \textcolor{blue}{Diodo~4} & \textcolor{blue}{Retificador (RQ2)} & \textcolor{blue}{Falha (diodo~4)} & \textcolor{blue}{Saudável} & \textcolor{blue}{Não confirmada\textsuperscript{a}} & \textcolor{blue}{Elétrica isolada} & \textcolor{blue}{736} \\
        \textcolor{blue}{Diodo~5} & \textcolor{blue}{Retificador (RQ2)} & \textcolor{blue}{Falha (diodo~5)} & \textcolor{blue}{Saudável} & \textcolor{blue}{Não confirmada\textsuperscript{a}} & \textcolor{blue}{Elétrica isolada} & \textcolor{blue}{479} \\
        \textcolor{blue}{Diodo~6} & \textcolor{blue}{Retificador (RQ2)} & \textcolor{blue}{Falha (diodo~6)} & \textcolor{blue}{Saudável} & \textcolor{blue}{Não confirmada\textsuperscript{a}} & \textcolor{blue}{Elétrica isolada} & \textcolor{blue}{658} \\
        \textcolor{blue}{Retificador completo} & \textcolor{blue}{Retificador (RQ2)} & \textcolor{blue}{Falha total} & \textcolor{blue}{Saudável} & \textcolor{blue}{Máxima do ramo} & \textcolor{blue}{Elétrica isolada} & \textcolor{blue}{736} \\
        \addlinespace
        \textcolor{blue}{Máquina balanceada (s/ parafusos)} & \textcolor{blue}{Moto-máquina (RQ1)} & \textcolor{blue}{Saudável} & \textcolor{blue}{Balanceada (linha de base)} & \textcolor{blue}{---} & \textcolor{blue}{Mecânica isolada} & \textcolor{blue}{1.450} \\
        \textcolor{blue}{Máquina desalinhada} & \textcolor{blue}{Moto-máquina (RQ1)} & \textcolor{blue}{Saudável} & \textcolor{blue}{Desalinhada} & \textcolor{blue}{Não caracterizada\textsuperscript{b}} & \textcolor{blue}{Mecânica isolada} & \textcolor{blue}{1.578} \\
        \textcolor{blue}{Máquina desbalanceada (RET~6~PARAF)} & \textcolor{blue}{Moto-máquina (RQ1)} & \textcolor{blue}{Saudável} & \textcolor{blue}{Desbalanceada} & \textcolor{blue}{Baixa (6 parafusos)} & \textcolor{blue}{Mecânica isolada} & \textcolor{blue}{1.453} \\
        \textcolor{blue}{Máquina desbalanceada (RET~14~PARAF)} & \textcolor{blue}{Moto-máquina (RQ1)} & \textcolor{blue}{Saudável} & \textcolor{blue}{Desbalanceada} & \textcolor{blue}{Intermediária (14 parafusos)} & \textcolor{blue}{Mecânica isolada} & \textcolor{blue}{1.450} \\
        \textcolor{blue}{Máquina desbalanceada, desalinhada (RET~15~PARAF)} & \textcolor{blue}{Moto-máquina (RQ1)} & \textcolor{blue}{Saudável} & \textcolor{blue}{Desbalanceada + desalinhada} & \textcolor{blue}{Alta (15 parafusos)} & \textcolor{blue}{Mecânica isolada} & \textcolor{blue}{1.443} \\
        \addlinespace
        \textcolor{blue}{Balanceada e alinhada + Diodo~2} & \textcolor{blue}{Moto-máquina + retificador\textsuperscript{c}} & \textcolor{blue}{Falha (diodo~2)} & \textcolor{blue}{Saudável} & \textcolor{blue}{Não confirmada\textsuperscript{a}} & \textcolor{blue}{Combinada} & \textcolor{blue}{1.572} \\
        \textcolor{blue}{Balanceada e alinhada + Diodo~3} & \textcolor{blue}{Moto-máquina + retificador\textsuperscript{c}} & \textcolor{blue}{Falha (diodo~3)} & \textcolor{blue}{Saudável} & \textcolor{blue}{Não confirmada\textsuperscript{a}} & \textcolor{blue}{Combinada} & \textcolor{blue}{1.589} \\
        \textcolor{blue}{Balanceada e alinhada + Diodo~6} & \textcolor{blue}{Moto-máquina + retificador\textsuperscript{c}} & \textcolor{blue}{Falha (diodo~6)} & \textcolor{blue}{Saudável} & \textcolor{blue}{Não confirmada\textsuperscript{a}} & \textcolor{blue}{Combinada} & \textcolor{blue}{1.506} \\
        \textcolor{blue}{Balanceada e alinhada + Retificador completo} & \textcolor{blue}{Moto-máquina + retificador\textsuperscript{c}} & \textcolor{blue}{Falha total} & \textcolor{blue}{Saudável} & \textcolor{blue}{Máxima do ramo} & \textcolor{blue}{Combinada} & \textcolor{blue}{1.440} \\
        \midrule
        \multicolumn{6}{l}{\textbf{\textcolor{blue}{Total}}} & \textbf{\textcolor{blue}{19.501}} \\
        \bottomrule
    \end{tabular}\\[0.3em]
    \raggedright\footnotesize\textcolor{blue}{\textsuperscript{a}~A numeração \texttt{Diodo~1}--\texttt{6} pode indicar a posição do diodo afetado no retificador trifásico ou uma contagem ordinal de diodos simultaneamente afetados; a semântica exata depende da topologia da bancada de retificador, ainda pendente de detalhamento (Seção~\ref{subsec:bancada_rq2}), por isso não é tratada aqui como escala de severidade confirmada. \textsuperscript{b}~O grau de desalinhamento (por exemplo, deslocamento angular/paralelo em mm ou graus) não foi quantificado nesta pesquisa, diferente do desbalanceamento (Seção~\ref{cap:limitacoes}). \textsuperscript{c}~As quatro condições combinadas exigem, por definição do rótulo, a máquina mecanicamente saudável operando simultaneamente à falha elétrica -- ou seja, envolvem necessariamente as duas bancadas (Seção~\ref{sec:bancadas}), e não apenas a bancada de retificador isoladamente; o arranjo físico exato dessa operação conjunta (mesmo ponto de acoplamento comum entre as duas bancadas) ainda depende do detalhamento pendente da topologia do retificador (Seção~\ref{subsec:bancada_rq2}).}
\end{table}

\begin{figure}[H]
    \centering
    \includegraphics[width=0.82\textwidth]{findings_labels}
    \caption{Distribuição de amostras pelas 17 condições experimentais -- de 479 a 1.589 amostras por condição.}
    \label{fig:res_labels}
\end{figure}

A assinatura harmônica de corrente ($H_{Ia}$, valor médio) para cinco condições representativas é apresentada na Figura~\ref{fig:res_signature}. A fundamental ($H_1 \approx 100\%$ em todas as condições) foi omitida por não carregar informação diagnóstica. O contraste visual é imediato: as condições de falha elétrica (\texttt{DIODO~3}, \texttt{RETIFICADOR COMPLETO}) apresentam amplitudes de 5 a 100 vezes maiores que as condições mecânicas ou saudável nas ordens harmônicas baixas (por exemplo, $H_2$ chega a 20\% para \texttt{DIODO~3}, contra 0,1--4\% para as demais condições mostradas), enquanto máquina balanceada, desalinhada e desbalanceada permanecem visualmente próximas entre si ao longo de quase todo o espectro mostrado -- um primeiro indício visual, confirmado estatisticamente na Seção~\ref{sec:resultados_sensibilidade}, de que a assinatura de falha mecânica é substancialmente mais sutil que a de falha elétrica.

\begin{figure}[H]
    \centering
    \includegraphics[width=0.85\textwidth]{findings_signature}
    \caption{Assinatura harmônica de corrente ($H_{Ia}$, média) para cinco condições representativas (H2--H25).}
    \label{fig:res_signature}
\end{figure}

Durante a validação da base de dados, um achado metodológico relevante foi documentado: as colunas \texttt{H00}--\texttt{H50~max} aparecem com nome idêntico nas duas planilhas de origem (\textsc{Harmônicas} e \textsc{Interharmônicas}, Seção~\ref{sec:bancadas}). A Figura~\ref{fig:res_hmax} compara diretamente os valores do canal \texttt{H05~Ia~max} entre as duas planilhas, pareados por rótulo e por proximidade temporal (tolerância de 90~s). A correlação global entre as duas séries é de 0,97 -- os pontos seguem visivelmente a diagonal $y=x$ por grupo de condição --, mas 0\% dos pares possui valor exatamente idêntico, e a comparação par-a-par completa (306 canais, 5.967.306 pares) revelou 85,81\% de divergência, com diferença absoluta máxima de 136,97 pontos percentuais. A conclusão adotada -- e mantida no espaço de atributos de todos os modelos deste capítulo -- é que não se trata de uma coluna duplicada, mas de duas fontes de medição distintas, preservadas como atributos separados.

\begin{figure}[H]
    \centering
    \includegraphics[width=0.72\textwidth]{findings_hmax_scatter}
    \caption{Comparação do canal H05~Ia~max entre a planilha Harmônicas e a planilha Interharmônicas (mesmo nome de coluna), pareado por rótulo e timestamp mais próximo.}
    \label{fig:res_hmax}
\end{figure}

Por fim, a Figura~\ref{fig:res_zero} caracteriza o piso de ruído do analisador de energia: o percentual de leituras exatamente zero cresce de forma aproximadamente monotônica a partir da ordem harmônica~20, ultrapassando 10\% a partir de $H_{30}$ e atingindo entre 20\% e 27\% nas ordens acima de $H_{36}$. Esse comportamento -- não uma característica física do sistema, mas uma limitação de resolução do instrumento nas ordens mais altas -- foi considerado na seleção de atributos empregada na desagregação de cargas (Seção~\ref{sec:resultados_rq2}) e ajuda a explicar por que os canais espectrais mais relevantes, tanto para diagnóstico (Seção~\ref{sec:resultados_rq1}) quanto para desagregação, concentram-se predominantemente em ordens baixas e médias.

\begin{figure}[H]
    \centering
    \includegraphics[width=0.85\textwidth]{findings_zero_pct}
    \caption{Percentual de valores exatamente zero por ordem harmônica (corrente Ia, média, todas as condições).}
    \label{fig:res_zero}
\end{figure}

\subsection{Diagnóstico de falha nas 17 condições --- evidência primária de RQ1}
\label{sec:resultados_rq1}

O mapa de assinatura da Figura~\ref{fig:res_heatmap} sintetiza, para as 17 condições e as ordens harmônicas $H_2$--$H_{25}$ de corrente Ia, o mesmo contraste já indicado na Figura~\ref{fig:res_signature}, agora em escala completa: as seis condições de falha elétrica (\texttt{DIODO~1}--\texttt{6}) dominam visualmente a escala de cor (compressão $\sqrt{\cdot}$ aplicada justamente para não obscurecer as demais condições), com amplitude máxima de até 50\% em $H_2$ (\texttt{DIODO~1}, \texttt{DIODO~4}). As \textcolor{blue}{dez} condições associadas à bancada moto-máquina permanecem visualmente escuras apenas em $H_5$ -- inclusive \texttt{RETIFICADOR COMPLETO}, que combina retificador com máquina balanceada e por isso herda parcialmente o padrão elétrico. Essa concentração de energia espectral nas ordens baixas, específica da falha elétrica, é a base física sobre a qual os classificadores desta seção operam.

\begin{figure}[H]
    \centering
    \includegraphics[width=\textwidth]{full_heatmap}
    \caption{Mapa de assinatura: amplitude média de corrente Ia (H2--H25) para as 17 condições experimentais, em escala de cor $\sqrt{\cdot}$.}
    \label{fig:res_heatmap}
\end{figure}

A importância de atributos via SHAP para um classificador \textit{Random Forest} treinado sobre as 17 classes (Figura~\ref{fig:res_shap_full}) confirma essa leitura: os 15 atributos mais relevantes são majoritariamente harmônicos de corrente de ordem baixa a média ($H_{19}\_I_c$, $H_{04}\_I_b$, $H_{05}\_I_b$, $H_{23}\_I_b$, entre outros), com a tensão aparecendo apenas na 10\textsuperscript{a} posição ($H_{05}\_V_a\_avg$). Esse predomínio de corrente sobre tensão é consistente com a premissa metodológica da MCSA clássica \citep{thomson2001,nandi2005} -- o defeito elétrico está fisicamente no caminho de corrente (retificador/diodo) -- e será contrastado diretamente, na Seção~\ref{sec:resultados_sensibilidade}, com o padrão inverso observado quando a influência elétrica é removida da análise.

\begin{figure}[H]
    \centering
    \includegraphics[width=0.78\textwidth]{full_shap}
    \caption{Importância SHAP (Random Forest, 17 classes) -- 15 atributos mais relevantes.}
    \label{fig:res_shap_full}
\end{figure}

\textcolor{blue}{Seis modelos clássicos de aprendizado supervisionado (Random Forest, XGBoost, SVM, MLP, k-NN, Regressão Logística -- Seção~\ref{subsec:modelos_classicos}) foram comparados sob o protocolo corrigido descrito na Seção~\ref{sec:protocolo}: divisão treino/teste cronológica 70/30 dentro de cada rótulo, com seleção de atributos (984 canais candidatos -- harmônicos e inter-harmônicos, Seção~\ref{sec:atributos} -- reduzidos a 40) reajustada apenas sobre as linhas de treino (Figura~\ref{fig:res_classical} e Tabela~\ref{tab:res_classical}). Sob esse protocolo, o MLP lidera com 90,45\% de acurácia e F1-macro de 0,908; Random Forest, k-NN e SVM formam um grupo intermediário (F1-macro entre 0,898 e 0,904); XGBoost e Regressão Logística ficam mais atrás (F1-macro de 0,883 e 0,860, respectivamente). Comparado à versão preliminar deste protocolo -- que reportava 92,29\% de acurácia para o SVM sob divisão aleatória estratificada e seleção de atributos ajustada na base inteira antes da divisão --, a acurácia do SVM sob o protocolo corrigido cai para 87,10\%: uma queda de cerca de 5 pontos percentuais, na direção e magnitude esperadas por uma validação que deixa de testar contra aquisições vizinhas da mesma sessão (Seção~\ref{sec:protocolo}).}

\begin{figure}[H]
    \centering
    \includegraphics[width=0.78\textwidth]{full_classical}
    \caption{\textcolor{blue}{Comparação de 6 modelos clássicos sob divisão temporal 70/30, seleção de atributos ajustada só no treino (17 classes, pool de 984 canais candidatos).}}
    \label{fig:res_classical}
\end{figure}

\begin{table}[H]
    \centering
    \caption{\textcolor{blue}{Modelos clássicos de referência, 17 classes, divisão treino/teste cronológica 70/30, seleção de atributos ajustada só no treino, pool de 984 canais candidatos.}}
    \label{tab:res_classical}
    \begin{tabular}{@{}lrr@{}}
        \toprule
        \textbf{Modelo} & \textbf{Acurácia} & \textbf{F1-macro} \\
        \midrule
        \textcolor{blue}{MLP} & \textcolor{blue}{90,45\%} & \textcolor{blue}{0,908} \\
        \textcolor{blue}{k-NN} & \textcolor{blue}{87,81\%} & \textcolor{blue}{0,905} \\
        \textcolor{blue}{Random Forest} & \textcolor{blue}{88,62\%} & \textcolor{blue}{0,904} \\
        \textcolor{blue}{SVM} & \textcolor{blue}{87,10\%} & \textcolor{blue}{0,898} \\
        \textcolor{blue}{XGBoost} & \textcolor{blue}{87,90\%} & \textcolor{blue}{0,883} \\
        \textcolor{blue}{Regressão Logística} & \textcolor{blue}{86,93\%} & \textcolor{blue}{0,860} \\
        \bottomrule
    \end{tabular}
\end{table}

\textcolor{blue}{Para as seis arquiteturas sequenciais introduzidas na Seção~\ref{subsec:rnns}, uma primeira rodada de comparação rápida (janela de 8 aquisições, 20 atributos, 6 épocas -- baseline, distinto do treinamento completo de 20 épocas detalhado na Seção~\ref{sec:resultados_rnn}) foi retreinada sob o protocolo corrigido: divisão temporal 75/25 por rótulo (Seção~\ref{sec:protocolo}) aplicada às aquisições \emph{antes} do janelamento, com as janelas de treino e de validação geradas separadamente em cada metade. Essa correção responde diretamente à observação do Pacheco de que, com tamanho de janela 8 e passo 4, janelas vizinhas compartilham metade de suas aquisições -- sob a divisão aleatória originalmente usada, uma fração das janelas de validação era quase idêntica a janelas de treino; uma checagem de amostragem confirmou zero janelas de validação idênticas a janelas de treino sob a divisão temporal. Sob esse protocolo mais honesto, a ordenação entre arquiteturas muda substancialmente e a acurácia geral cai (Figura~\ref{fig:res_dl}): GRU lidera com 67,8\% de acurácia de validação, seguido por SimpleRNN (66,0\%), LSTM (64,5\%), BiLSTM (64,3\%), CNN+LSTM (63,4\%) e Transformer (62,6\%) -- uma queda de 10 a 20 pontos percentuais frente à versão com vazamento, e uma inversão quase completa da ordenação anterior, que apontava o Transformer como líder. Esse resultado, obtido sob uma única semente de inicialização por arquitetura (ver limitação registrada na Seção~\ref{cap:limitacoes}), deve ser lido como uma ordenação observada nesta execução específica, não como diferença estatisticamente estabelecida entre arquiteturas adjacentes; mas a magnitude da queda geral de acurácia é, em si, a evidência mais direta reunida nesta pesquisa de que a divisão aleatória sobre janelas sobrepostas inflava artificialmente o desempenho reportado das redes sequenciais.}

\begin{figure}[H]
    \centering
    \includegraphics[width=0.78\textwidth]{full_dl}
    \caption{\textcolor{blue}{Comparação rápida de 6 arquiteturas de rede neural com contexto temporal, sob divisão temporal (janela de 8 aquisições, 20 atributos, 6 épocas).}}
    \label{fig:res_dl}
\end{figure}

Como evidência complementar ao diagnóstico -- não uma terceira saída independente --, um modelo BiLSTM foi treinado para prever a assinatura espectral um passo à frente (canal $H_{02}\_I_a\_avg$). A Figura~\ref{fig:res_forecast} mostra a dispersão entre valor real e valor previsto no conjunto de validação: os pontos seguem estreitamente a diagonal $y=x$ em toda a faixa de amplitude observada (de próximo de zero a mais de 60\%), com correlação de 0,996 entre série real e prevista. \textcolor{blue}{Essa correlação alta, por si só, não demonstra que o modelo aprendeu algo além da forte autocorrelação temporal do próprio sinal: um preditor de persistência trivial ($\hat{y}_{t+1}=y_t$, prever que o próximo valor é igual ao atual), avaliado no mesmo conjunto de validação e no mesmo canal, atinge correlação de 0,9959 -- ligeiramente \textbf{acima} da do BiLSTM (0,9952) -- e MAE de 0,387, menos da metade do MAE do BiLSTM (0,777). Ou seja, sob esse teste direto, o preditor de persistência iguala ou supera o BiLSTM nas duas métricas: a correlação de 0,996 não é evidência de que a rede aprendeu uma dinâmica útil além de "o próximo valor é parecido com o atual". Prever um canal espectral um passo à frente não demonstra, portanto, previsão de falha nem detecção precoce de anomalia -- a linguagem apropriada para este resultado é diagnóstico de condição a partir do sinal já instalado, com previsão de trajetória futura permanecendo como objetivo em aberto até que um experimento longitudinal de degradação (Seção~\ref{cap:limitacoes}) permita comparar o BiLSTM contra a persistência em uma janela onde as duas realmente divergem.}

\begin{figure}[H]
    \centering
    \includegraphics[width=0.68\textwidth]{full_forecast}
    \caption{Previsão da assinatura futura (H02\_Ia\_avg, um passo à frente) por um modelo BiLSTM: valor real \textit{vs.} previsto no conjunto de validação.}
    \label{fig:res_forecast}
\end{figure}

Uma primeira rodada de desagregação de cargas, restrita a atributos harmônicos, é apresentada na Figura~\ref{fig:res_disagg_prelim} por completude histórica. Como detalhado na Seção~\ref{sec:resultados_rq2}, o $R^2$ calculado por dobra individual sob validação \textit{leave-one-condition-out} é estruturalmente degenerado -- o alvo é constante dentro de cada dobra, de modo que o $R^2$ do \texttt{scikit-learn} recai no caso especial de variância nula tanto para o alvo ``máquina'' quanto para o alvo ``fonte elétrica''. Essa figura não deve, portanto, ser lida como evidência quantitativa válida de RQ2; a evidência primária correta é apresentada na Seção~\ref{sec:resultados_rq2} (Figura~\ref{fig:res_nilm}).

\begin{figure}[H]
    \centering
    \includegraphics[width=0.8\textwidth]{full_disagg}
    \caption{Primeira rodada de desagregação de cargas (apenas harmônicas): $R^2$ por dobra de validação \textit{leave-one-condition-out}, por alvo (máquina \textit{vs.} fonte elétrica). Mantida por completude histórica -- ver ressalva metodológica na Seção~\ref{sec:resultados_rq2}.}
    \label{fig:res_disagg_prelim}
\end{figure}

\textcolor{blue}{Uma síntese dos principais achados desta seção em linguagem acessível, sem jargão estatístico -- útil como material de apoio a uma apresentação de qualificação para público não especializado, mas que reapresenta dados já discutidos nesta e na próxima seção, sem constituir evidência adicional -- foi deslocada do núcleo de Resultados para o Apêndice~C, atendendo à observação do orientador de que essa síntese não deveria competir, na leitura do capítulo, com a evidência quantitativa primária.}

\textcolor{blue}{
\subsection{Classificação hierárquica: quatro macro-categorias}
\label{sec:resultados_macro}
Como segunda das três avaliações separadas de RQ1 solicitadas pelo orientador -- classificação global das 17 condições (Seção~\ref{sec:resultados_rq1}), classificação hierárquica das quatro categorias principais (esta seção) e diagnóstico exclusivamente mecânico (Seção~\ref{sec:resultados_sensibilidade}) --, os seis modelos clássicos foram retreinados sobre o mesmo pool de 984 canais candidatos e o mesmo protocolo de divisão temporal (Seção~\ref{sec:protocolo}), mas com o rótulo de 17 classes substituído pela categoria científica de primeiro nível da Tabela~\ref{tab:res_taxonomia_hierarquica} (saudável, elétrica isolada, mecânica isolada, combinada). Diferente da reorganização puramente interpretativa descrita na Seção~\ref{sec:resultados_bd}, esta é uma reavaliação computacional de fato: os modelos são retreinados do zero sobre o rótulo agregado, não uma reagregação post-hoc das previsões do classificador de 17 classes.

\begin{table}[H]
    \centering
    \caption{Classificação das 4 macro-categorias, divisão temporal 70/30, pool de 984 canais candidatos.}
    \label{tab:res_macro}
    \begin{tabular}{@{}lrr@{}}
        \toprule
        \textbf{Modelo} & \textbf{Acurácia} & \textbf{F1-macro} \\
        \midrule
        Random Forest & 97,97\% & 0,954 \\
        XGBoost & 97,28\% & 0,942 \\
        k-NN & 95,71\% & 0,925 \\
        MLP & 91,68\% & 0,878 \\
        SVM & 87,18\% & 0,817 \\
        Regressão Logística & 83,42\% & 0,784 \\
        \bottomrule
    \end{tabular}
\end{table}

A acurácia é substancialmente mais alta do que a observada nas 17 classes (Tabela~\ref{tab:res_classical}, 87--90\%) -- resultado esperado, já que separar apenas quatro categorias amplas (saudável \textit{vs.} elétrica \textit{vs.} mecânica \textit{vs.} combinada) é um problema estritamente mais fácil do que distinguir os 17 subtipos dentro delas, e confirma que a dificuldade da classificação de 17 classes está concentrada na distinção de subtipo e severidade dentro de cada ramo -- não na distinção entre os quatro ramos em si, que a assinatura espectral separa com desempenho próximo de 98\% de acurácia (Random Forest).

Sob validação cruzada temporal em blocos de 5 dobras, porém, essa mesma tarefa exibe uma dispersão entre dobras muito maior do que a observada nas 17 classes ou no subconjunto mecânico: entre 14,0 (SVM) e 17,7 (k-NN) pontos percentuais de desvio-padrão de F1-macro, contra 3 a 4 pontos percentuais nas 17 classes (Tabela~\ref{tab:res_cv}). Essa dispersão elevada não deve ser lida como instabilidade do pipeline, mas como consequência direta de reduzir o problema a apenas quatro classes: com poucas classes, o F1-macro de uma dobra é dominado pelo desempenho na classe de menor representação dentro daquela dobra especificamente -- notavelmente a categoria \texttt{saudável}, com apenas 1.440 amostras no total (uma única condição experimental, Tabela~\ref{tab:res_taxonomia_hierarquica}) distribuídas em blocos temporais de cerca de 288 amostras por dobra --, de modo que uma única dobra em que essa classe é mal representada ou mal prevista derruba desproporcionalmente a média macro daquela dobra. Esse é o mesmo efeito, em escala ampliada, já discutido para a dispersão das 17 classes e do subconjunto mecânico (Seção~\ref{sec:resultados_robustez}): dobras temporais em blocos não são permutáveis entre si por construção, e a redução do número de classes torna essa não-permutabilidade mais visível na métrica agregada, não menos real.
}

\subsection{Sensibilidade mecânica isolada --- apoio interpretativo a RQ1}
\label{sec:resultados_sensibilidade}

Esta seção não constitui um objetivo autônomo da pesquisa, mas uma sub-análise de controle metodológico: restringindo o universo de análise às seis condições mecânicas puras (Tabela~\ref{tab:res_taxonomia}), sem nenhuma falha elétrica presente, isola-se a variável mecânica do efeito de confusão elétrico -- qualquer harmônico de corrente é, por construção, sensível à presença de uma carga não linear, o que poderia mascarar a real capacidade diagnóstica sobre falhas mecânicas se avaliada apenas dentro do conjunto completo de 17 classes.

Foram testados por ANOVA de um fator 606 canais harmônicos e 612 inter-harmônicos, com correção de Benjamini-Hochberg (FDR, $\alpha=0{,}05$) sobre o conjunto completo testado simultaneamente: 97,0\% dos canais harmônicos e 99,3\% dos inter-harmônicos permaneceram significativos após a correção. \textcolor{blue}{Esses 606+612=1.218 canais não são o mesmo universo dos 984 canais candidatos citados na Seção~\ref{sec:atributos}: aqui a varredura ANOVA é feita diretamente sobre as tabelas em formato longo (\textit{long format}) de cada planilha de origem, contando toda combinação família/variável/fase/ordem/estatística presente nelas, e restrita às seis condições mecânicas; os 984 canais candidatos, por outro lado, são colunas já pivotadas da matriz combinada (Seção~\ref{sec:bd}), deduplicadas pela família espectral (\texttt{parse\_feature\_key}) e cobrindo as 17 condições. São duas contagens de coisas diferentes, não uma inconsistência a reconciliar -- e é justamente por não terem sido diferenciadas explicitamente no texto que pareciam contraditórias.} \textcolor{blue}{Com mais de 1.400 amostras por condição mecânica, esse resultado reflete majoritariamente o elevado poder estatístico do desenho experimental, não necessariamente relevância prática de cada canal individual -- e a causa desse poder estatístico exige uma ressalva adicional, não apenas o tamanho nominal da amostra: como discutido na Seção~\ref{sec:protocolo}, cada condição corresponde a uma única sessão de aquisição contínua, de modo que as 1.400+ aquisições de uma mesma condição não são unidades experimentais independentes, mas observações sucessivas e temporalmente autocorrelacionadas da mesma execução -- a mesma pseudorreplicação que motiva a divisão temporal na validação dos classificadores (Seção~\ref{sec:protocolo}) também infla o poder estatístico nominal do teste ANOVA aqui, e a correção de Benjamini-Hochberg (FDR) corrige apenas a multiplicidade de testes, não essa falta de independência entre observações. Por isso a interpretação desta seção apoia-se em $\eta^2$ (tamanho de efeito multi-grupo) e no delta de Cliff (contrastes par-a-par) em vez do valor-$p$ isoladamente -- $\eta^2$ e o delta de Cliff não escalam artificialmente com o número de aquisições da forma que o valor-$p$ escala, o que os torna leitura mais robusta da magnitude do efeito sob pseudorreplicação, ainda que também não sejam imunes a viés proveniente de autocorrelação temporal dentro de uma mesma sessão; a taxa de significância de 97,0\%/99,3\% relatada acima, especificamente, deve ser lida com essa ressalva, e não como evidência de que quase todo o espaço espectral é diagnosticamente relevante.}

A Figura~\ref{fig:res_eta} mostra os oito canais harmônicos e os seis inter-harmônicos de maior $\eta^2$. O achado mais marcante desta sub-análise está do lado inter-harmônico: as componentes de ordem 0 e 1 (aproximadamente contínua) da corrente concentram os maiores tamanhos de efeito de todo o levantamento, com $\eta^2$ entre 0,90 e 0,95 para $IH_0\_I_c$, $IH_0\_I_b$, $IH_0\_I_a$, $IH_1\_I_b$, $IH_1\_I_c$ e $IH_1\_I_a$ -- todos acima do maior $\eta^2$ harmônico observado (0,863, em $H_{10}\_I_b\_max$). Este resultado converge diretamente com a motivação teórica estabelecida no Capítulo~\ref{cap:fundamentacao} a partir do achado de \citet{efe2016} de que componentes inter-harmônicas podem carregar distorção percentual superior à da harmônica adjacente: aqui, para o problema de sensibilidade mecânica isolada, as inter-harmônicas não apenas igualam, mas superam sistematicamente as harmônicas como atributo discriminante.

\begin{figure}[H]
    \centering
    \includegraphics[width=\textwidth]{tecnico_eta}
    \caption{Top~8 canais harmônicos e top~6 inter-harmônicos por $\eta^2$ (6 condições mecânicas).}
    \label{fig:res_eta}
\end{figure}

A importância SHAP de um classificador \textit{Random Forest} restrito às seis condições mecânicas (Figura~\ref{fig:res_shap_mech}) revela um padrão que se inverte em relação ao observado no conjunto completo de 17 classes (Figura~\ref{fig:res_shap_full}): o atributo mais importante é de tensão ($H_{05}\_V_a\_avg$), e cinco dos quinze atributos mais relevantes são de tensão -- contra apenas um de quinze no conjunto completo (Figura~\ref{fig:res_leigos_volt}). \textcolor{blue}{Essa inversão tensão/corrente sugere maior conteúdo discriminativo em determinados atributos de tensão no subconjunto mecânico. Uma leitura fisicamente plausível é que desbalanceamento e desalinhamento afetam primariamente o entreferro e a força eletromotriz induzida, refletindo-se mais na tensão -- mas essa interpretação permanece hipótese, não fato demonstrado por esta pesquisa: a tensão medida no ponto de acoplamento comum depende também da impedância da fonte, da rede, dos cabos e da instrumentação, de modo que o mecanismo eletromecânico associado precisa ser investigado considerando a topologia e a impedância da alimentação} -- uma leitura que será retomada com mais profundidade, à luz da literatura de MCSA \citep{thomson2001,nandi2005,benamira2022}, no Capítulo~\ref{cap:discussao}.

\begin{figure}[H]
    \centering
    \includegraphics[width=0.78\textwidth]{tecnico_shap}
    \caption{Importância SHAP (Random Forest, condições mecânicas) -- 15 atributos mais relevantes, coloridos por tensão (V) ou corrente (I).}
    \label{fig:res_shap_mech}
\end{figure}

\textcolor{blue}{Sob o protocolo corrigido (Seção~\ref{sec:protocolo}), o mesmo pool de 984 canais foi reavaliado restrito às seis condições mecânicas (Tabela~\ref{tab:res_classical_mech}): SVM e Random Forest lideram, com acurácia de 82,90\%/82,98\% e F1-macro de 0,827/0,826. Comparado ao SVM nas 17 classes (F1-macro 0,898, Tabela~\ref{tab:res_classical}), isolar a mecânica ainda derruba o desempenho -- queda de cerca de 7 pontos percentuais de F1-macro --, mas a queda é menor do que a reportada pela versão preliminar deste protocolo (13 pontos, de 0,939/0,937 para 0,806/0,801): parte daquela diferença maior era o próprio SVM nas 17 classes estar artificialmente inflado pelo vazamento de seleção de atributos (Seção~\ref{sec:resultados_rq1}), não uma queda real de 13 pontos ao isolar a mecânica. A maior sutileza da assinatura mecânica isolada, já sugerida qualitativamente pela Figura~\ref{fig:res_signature}, permanece real -- só que menor do que o número antigo indicava.}

\begin{table}[H]
    \centering
    \caption{\textcolor{blue}{Modelos clássicos de referência, 6 classes mecânicas puras, divisão treino/teste cronológica 70/30, pool de 984 canais candidatos.}}
    \label{tab:res_classical_mech}
    \begin{tabular}{@{}lrr@{}}
        \toprule
        \textbf{Modelo} & \textbf{Acurácia} & \textbf{F1-macro} \\
        \midrule
        \textcolor{blue}{SVM} & \textcolor{blue}{82,90\%} & \textcolor{blue}{0,827} \\
        \textcolor{blue}{Random Forest} & \textcolor{blue}{82,98\%} & \textcolor{blue}{0,826} \\
        \textcolor{blue}{MLP} & \textcolor{blue}{82,30\%} & \textcolor{blue}{0,811} \\
        \textcolor{blue}{k-NN} & \textcolor{blue}{80,75\%} & \textcolor{blue}{0,798} \\
        \textcolor{blue}{Regressão Logística} & \textcolor{blue}{80,98\%} & \textcolor{blue}{0,770} \\
        \textcolor{blue}{XGBoost} & \textcolor{blue}{77,50\%} & \textcolor{blue}{0,718} \\
        \bottomrule
    \end{tabular}
\end{table}

\begin{figure}[H]
    \centering
    \includegraphics[width=0.7\textwidth]{tecnico_models}
    \caption{\textcolor{blue}{F1-macro por modelo clássico sob divisão temporal, problema de 17 classes \textit{vs.} subconjunto de 6 classes mecânicas puras.}}
    \label{fig:res_models_mech}
\end{figure}

Contrastes par-a-par (delta de Cliff, cada condição \textit{vs.} baseline saudável) revelam heterogeneidade qualitativa entre tipos de falha mecânica, não apenas graus diferentes de uma mesma assinatura: desalinhamento separa-se do baseline predominantemente em canais de \textbf{tensão} ($H_5\_V_b\_avg/max$, $\delta=1{,}000$; $H_{21}\_V_b\_avg$, $\delta=-0{,}9999$), enquanto desbalanceamento (\texttt{RET~6~PARAF}) separa-se em \textbf{corrente} de ordem~9 ($\delta=0{,}816$) e tensão de ordem~27 ($\delta=-0{,}798$). Essa heterogeneidade indica que diferentes modos de falha mecânica produzem assinaturas espectrais qualitativamente distintas -- achado relevante para o desenho de um espaço de atributos diagnóstico que não trate ``falha mecânica'' como uma categoria espectralmente homogênea.

\subsection{Desagregação de cargas --- evidência primária de RQ2}
\label{sec:resultados_rq2}

Do espaço de 984 canais espectrais candidatos (666 harmônicos + 318 inter-harmônicos, Seção~\ref{sec:atributos}), 60 foram selecionados por F-score ANOVA com eliminação de redundância por correlação (limiar 0,95). \textcolor{blue}{Sob o protocolo corrigido (Seção~\ref{sec:protocolo}), essa seleção é ajustada excluindo inteiramente as quatro condições combinadas -- não apenas a retida em cada dobra --, resultando em 28 harmônicos, 29 inter-harmônicos, 1 de THD e 2 de THDe} -- uma distribuição praticamente balanceada entre harmônicas e inter-harmônicas, que já antecipa, antes de qualquer resultado de desempenho, que a etapa de seleção de atributos não descartou sistematicamente nenhum dos dois domínios.

Uma ressalva metodológica central precisa ser registrada antes da evidência quantitativa de RQ2, pois molda diretamente qual figura deste capítulo deve ser lida como resultado válido. A validação \textit{leave-one-condition-out} produz, por desenho, exatamente uma condição de teste por dobra -- logo o alvo (assinatura da máquina \emph{ou} da fonte elétrica) é constante dentro de cada dobra. Nessa condição, o $R^2$ do \texttt{scikit-learn} recai no caso especial de variância nula (0,0 ou 1,0 por definição, não uma medida de ajuste real) para os \textbf{dois} alvos -- não apenas para o alvo ``máquina'', como uma versão preliminar deste relatório havia indicado (Figura~\ref{fig:res_disagg_prelim}). O MAE permanece válido nessa configuração, por não depender da variância do alvo; \textcolor{blue}{o $R^2$ por canal, calculado sobre as previsões das quatro dobras agrupadas antes da pontuação, oferece uma estimativa exploratória de generalização entre as quatro condições combinadas disponíveis -- não uma medida com poder de generalização pleno, dado o número reduzido de condições independentes} -- procedimento que restaura variância real ao alvo, já que cada condição combinada contribui uma fonte elétrica diferente.

\textcolor{blue}{A Figura~\ref{fig:res_nilm} apresenta essa evidência corrigida: $R^2$ por canal, alvo igual à assinatura da fonte elétrica, para os 10 melhores canais entre os 60 avaliados, sob a seleção sem vazamento descrita acima. O melhor canal individual atinge $R^2=0{,}158$ ($IH_{46}\_I_c\_avg$); entre os quatro melhores, dois são inter-harmônicos ($IH_{46}\_I_c$, $IH_{48}\_I_c$, $R^2$ entre 0,142 e 0,158) e os outros dois são $THDe\_I_c\_avg$ e $H_{12}\_I_b\_avg$ ($R^2$ entre 0,105 e 0,122); entre os vinte melhores canais, catorze são inter-harmônicos e seis são harmônicos -- a mesma proporção aproximada já registrada na versão preliminar deste protocolo (treze de vinte), agora sob uma seleção que nunca viu as linhas das quatro condições combinadas. Esse predomínio de inter-harmônicas entre os canais mais recuperáveis é contraintuitivo frente à literatura de NILM revisada no Capítulo~\ref{cap:fundamentacao}, tradicionalmente concentrada em conteúdo harmônico inteiro \citep{djordjevic2017,bouhouras2017,li2019moea} -- ainda que consistente, mais uma vez, com a motivação original de \citet{efe2016} para investigar interharmônicas neste espaço de atributos. Dos 60 canais avaliados, 52 (86,7\%) apresentam $R^2$ negativo (contra 45/60=75\% na versão com vazamento) -- a seleção sem vazamento é mais honesta, não mais otimista: a desagregação continua seletiva por canal, não uniforme, e deve ser lida como evidência real, porém parcial, de viabilidade, mais ainda modesta agora do que antes.}

\begin{figure}[H]
    \centering
    \includegraphics[width=0.85\textwidth]{tecnico_nilm}
    \caption{\textcolor{blue}{$R^2$ por canal sob seleção de atributos sem vazamento (previsões agrupadas das 4 dobras), alvo = assinatura da fonte elétrica, 10 melhores de 60 canais avaliados.}}
    \label{fig:res_nilm}
\end{figure}

\textcolor{blue}{
\subsection{Ablação controlada: contribuição incremental de H e IH (RQ1 e RQ2)}
\label{sec:resultados_ablacao}

As seções anteriores usam ranking de atributos (SHAP, $\eta^2$, posição entre os canais mais bem recuperados) como evidência de que as componentes inter-harmônicas contribuem para o diagnóstico e para a desagregação -- um indício, não uma demonstração causal de ganho preditivo. Esta seção reporta o experimento decisivo: sob o mesmo particionamento e o mesmo modelo, comparar diretamente o desempenho usando apenas harmônicas (H), apenas inter-harmônicas (IH) e as duas combinadas (H+IH).

Para RQ1, o classificador SVM (Seção~\ref{sec:resultados_rq1}) foi reavaliado sob validação cruzada temporal em blocos de 5 dobras (Seção~\ref{sec:protocolo}), com o mesmo procedimento de seleção de atributos aplicado a cada um dos três pools de candidatos (H: 666 canais; IH: 318 canais; H+IH: 984 canais), gerando três séries de F1-macro pareadas dobra a dobra (Tabela~\ref{tab:res_ablacao}). Nas 17 classes, H e IH sozinhos ficam praticamente empatados (F1-macro médio 0,889 e 0,887), e H+IH ganha pouco sobre H isolado ($\Delta=+0{,}013$; teste de Wilcoxon pareado, $p=0{,}44$, não significativo com $n=5$ dobras) -- consistente com a falha elétrica já saturar o problema de 17 classes independentemente da origem espectral do atributo (Seção~\ref{sec:discussao_rq1}). Nas seis classes mecânicas puras, porém, o quadro muda por completo: IH sozinho (F1-macro 0,791) supera H sozinho (F1-macro 0,683) por quase 11 pontos percentuais, e H+IH (F1-macro 0,798) mal melhora sobre IH sozinho -- o ganho incremental de adicionar H a IH é pequeno; é IH que carrega a informação diagnóstica para falha mecânica, não H. $\Delta(\text{H+IH}-\text{H})=+0{,}116$ (Wilcoxon $p=0{,}063$): não atinge o limiar convencional de significância com apenas 5 dobras, mas é uma tendência forte e consistente com a análise de sensibilidade ($\eta^2$, Seção~\ref{sec:resultados_sensibilidade}), que já apontava as componentes de ordem 0-1 da corrente como as de maior tamanho de efeito de todo o levantamento.

\begin{table}[H]
    \centering
    \caption{Ablação H/IH/H+IH, SVM, F1-macro médio $\pm$ desvio padrão sobre 5 dobras temporais em blocos.}
    \label{tab:res_ablacao}
    \begin{tabular}{@{}lrrr@{}}
        \toprule
        \textbf{Problema} & \textbf{H} & \textbf{IH} & \textbf{H+IH} \\
        \midrule
        17 classes & 0,889 $\pm$ 0,038 & 0,887 $\pm$ 0,017 & 0,902 $\pm$ 0,033 \\
        6 classes (mecânica) & 0,683 $\pm$ 0,113 & 0,791 $\pm$ 0,050 & 0,798 $\pm$ 0,059 \\
        \bottomrule
    \end{tabular}
\end{table}

Para RQ2, a mesma lógica foi aplicada à desagregação (Seção~\ref{sec:resultados_rq2}): a seleção de atributos sem vazamento (Seção~\ref{sec:protocolo}) foi repetida três vezes, uma por pool de candidatos, escolhendo até 60 canais em cada uma, e o $R^2$ por canal (previsões agrupadas das 4 dobras \textit{leave-one-condition-out}, alvo = fonte elétrica) foi recalculado para cada pool (Tabela~\ref{tab:res_ablacao_rq2}). O resultado é o mais direto de toda esta pesquisa a favor da contribuição das inter-harmônicas: usando somente harmônicas, \textbf{nenhum} dos 60 canais selecionados atinge $R^2$ positivo (melhor canal: $R^2=-0{,}002$, i.e., não melhor do que prever a média) -- harmônicas isoladas não sustentam desagregação alguma neste espaço de atributos. Usando somente inter-harmônicas, o melhor canal atinge $R^2=0{,}880$, o maior valor observado em toda a análise de RQ2, embora 80\% dos 60 canais selecionados ainda fiquem negativos. Combinando H+IH, o melhor canal cai para $R^2=0{,}158$ (a mesma seleção reportada na Tabela~\ref{tab:res_nilm}) -- abaixo do pool só-IH, não acima dele.

\begin{table}[H]
    \centering
    \caption{Ablação H/IH/H+IH, RQ2, $R^2$ por canal (4 dobras agrupadas, alvo = fonte elétrica), até 60 canais selecionados por pool.}
    \label{tab:res_ablacao_rq2}
    \begin{tabular}{@{}lrrrr@{}}
        \toprule
        \textbf{Pool} & \textbf{Melhor $R^2$} & \textbf{$R^2$ médio} & \textbf{$R^2$ mediano} & \textbf{\% canais negativos} \\
        \midrule
        H & $-0{,}002$ & $-1{,}138$ & $-0{,}823$ & 100,0\% \\
        IH & $0{,}880$ & $-0{,}851$ & $-0{,}857$ & 80,0\% \\
        H+IH & $0{,}158$ & $-0{,}813$ & $-0{,}626$ & 86,7\% \\
        \bottomrule
    \end{tabular}
\end{table}

Esse padrão -- IH isolado superando H+IH combinado -- parece contraintuitivo à primeira vista, mas revela uma limitação metodológica real, não um erro: a seleção de atributos usa F-score ANOVA contra o rótulo de 17 classes (Seção~\ref{sec:atributos}) como critério, um \textit{proxy} para separabilidade de classe, não o alvo de regressão da desagregação em si. Quando H e IH competem na mesma classificação por F-score (pool H+IH), canais harmônicos com F-score mais alto para separar rótulos -- mas irrelevantes para reconstruir a assinatura da fonte elétrica -- ocupam vagas no corte de 60 que, no pool só-IH, sobrariam para o canal inter-harmônico de maior $R^2$ de regressão. Isso não invalida o achado central -- harmônicas isoladas nunca produzem um canal com $R^2$ positivo, em nenhum dos três pools --, mas indica que um critério de seleção alinhado ao objetivo de regressão (em vez de um \textit{proxy} de classificação) é um passo metodológico natural para a versão de artigo desta pesquisa.
}

\subsection{Redes neurais recorrentes --- curvas de treinamento}
\label{sec:resultados_rnn}

\textcolor{blue}{As seis arquiteturas sequenciais especificadas na Tabela~\ref{tab:res_rnn_spec} foram retreinadas por até 20 épocas (parada antecipada com paciência de 5, monitorando perda de validação) -- contra as 6 épocas da comparação rápida da Seção~\ref{sec:resultados_rq1} (Figura~\ref{fig:res_dl}) --, especificamente para expor a curva completa de convergência, mantendo a mesma divisão temporal treino/validação (Seção~\ref{sec:protocolo}) e os mesmos 20 atributos. Uma primeira execução, sob uma única semente de inicialização por arquitetura, produziu as curvas das Figuras~\ref{fig:res_rnn_curves}--\ref{fig:res_rnn_loss}; como essa execução isolada mostrou-se, na prática, sensível à semente utilizada (ver adiante), o treinamento foi repetido sob cinco sementes independentes por arquitetura, cujo resumo (Tabela~\ref{tab:res_rnn_multiseed}) é a evidência tratada como principal nesta seção e na Discussão.}

\begin{table}[H]
    \centering
    \caption{Especificação das seis arquiteturas de rede neural sequencial (Seção~\ref{subsec:rnns}).}
    \label{tab:res_rnn_spec}
    \begin{tabular}{@{}llr@{}}
        \toprule
        \textbf{Arquitetura} & \textbf{Camadas} & \textbf{Parâmetros} \\
        \midrule
        Transformer & Proj(32) $\to$ MHA(4$\times$32) $\to$ FF(64) $\to$ Pool $\to$ Dense(17) & 22.353 \\
        CNN+LSTM & Conv1D(32) $\to$ LSTM(64) $\to$ Dense(17) & 27.889 \\
        BiLSTM & BiLSTM(64$\times$2) $\to$ Dense(17) & 45.713 \\
        LSTM & LSTM(64) $\to$ Dense(17) & 22.865 \\
        GRU & GRU(64) $\to$ Dense(17) & 17.617 \\
        SimpleRNN & SimpleRNN(64) $\to$ Dense(17) & 6.545 \\
        \bottomrule
    \end{tabular}
\end{table}

\textcolor{blue}{A Figura~\ref{fig:res_rnn_curves} mostra acurácia de treino e de validação por época, individualmente para cada arquitetura, para a execução de semente única. Diferente da versão preliminar deste treinamento -- que reportava \textit{gap} negativo (validação acima de treino) em três das seis arquiteturas, atribuído ao efeito do \textit{dropout} --, sob divisão temporal todas as seis arquiteturas exibem \textit{gap} treino/validação positivo (acurácia de treino acima da de validação) em toda execução observada, seed única ou multi-semente. Esse padrão é o esperado para sobre-ajuste real, não um artefato do \textit{dropout}: sob a divisão aleatória anterior, o treino convergia sobre janelas quase idênticas às de validação, mascarando o sobre-ajuste; a divisão temporal, ao eliminar essa sobreposição, revela a diferença de desempenho que já existia. Parada antecipada interrompeu o treinamento tipicamente entre a 7ª e a 16ª época, a depender da arquitetura e da semente (paciência de 5 sobre a perda de validação), sem que praticamente nenhuma execução tenha alcançado as 20 épocas completas.}

\begin{figure}[H]
    \centering
    \includegraphics[width=\textwidth]{rnn_curves}
    \caption{\textcolor{blue}{Acurácia de treino (linha sólida) e validação (linha tracejada) por época sob divisão temporal, para cada uma das seis arquiteturas (execução de semente única).}}
    \label{fig:res_rnn_curves}
\end{figure}

\textcolor{blue}{A Figura~\ref{fig:res_rnn_loss} sobrepõe a perda de validação das seis arquiteturas ao longo do treinamento, para a execução de semente única. Nessa execução isolada, CNN+LSTM atingiu a maior acurácia de validação (67,0\%) e o Transformer a menor (64,3\%), com o maior \textit{gap} treino/validação das seis (24,0 pontos percentuais). \textbf{Esse resultado, no entanto, não se sustentou sob repetição}: como o próprio texto já alertava, uma execução de semente única não permite distinguir ordenação real de ruído de inicialização, e a repetição sob cinco sementes independentes (Tabela~\ref{tab:res_rnn_multiseed}) inverteu justamente o extremo mais chamativo dessa conclusão -- o Transformer, que na execução isolada aparecia como pior arquitetura, apresenta na verdade a maior acurácia média de validação entre as seis (68,6\% $\pm$ 2,0 pontos percentuais), e o CNN+LSTM, que liderava na execução isolada, apresenta a menor média (65,7\% $\pm$ 1,3 pontos percentuais). A ordenação média completa é Transformer (68,6\%) > LSTM (67,2\%) > BiLSTM (66,8\%) > SimpleRNN (66,4\%) > GRU (66,2\%) > CNN+LSTM (65,7\%), com desvios-padrão de 0,6 a 3,0 pontos percentuais entre sementes -- intervalos que se sobrepõem substancialmente entre arquiteturas adjacentes, de modo que a diferença entre a maioria dos pares não deve ser lida como estabelecida além do acaso. O padrão de \textit{gap} treino/validação, por sua vez, é mais estável entre sementes do que a acurácia: o Transformer mantém o maior \textit{gap} médio em todas as cinco execuções (20,6 $\pm$ 1,1 pontos percentuais), contra 9,8--16,2 pontos percentuais nas demais arquiteturas -- ou seja, apesar de sobre-ajustar mais consistentemente do que as demais, o Transformer generaliza, em média, tão bem ou melhor do que elas em termos absolutos, um padrão que só a repetição multi-semente tornou visível.}

\begin{table}[H]
    \centering
    \caption{\textcolor{blue}{Acurácia de validação e \textit{gap} treino/validação ao final do treinamento (média $\pm$ desvio padrão sobre 5 sementes independentes por arquitetura), sob divisão temporal.}}
    \label{tab:res_rnn_multiseed}
    \begin{tabular}{@{}lrr@{}}
        \toprule
        \textbf{Arquitetura} & \textbf{Acurácia de validação} & \textbf{\textit{Gap} treino/validação} \\
        \midrule
        \textcolor{blue}{Transformer} & \textcolor{blue}{68,6\% $\pm$ 2,0 pp} & \textcolor{blue}{20,6 $\pm$ 1,1 pp} \\
        \textcolor{blue}{LSTM} & \textcolor{blue}{67,2\% $\pm$ 0,9 pp} & \textcolor{blue}{11,0 $\pm$ 1,1 pp} \\
        \textcolor{blue}{BiLSTM} & \textcolor{blue}{66,8\% $\pm$ 2,0 pp} & \textcolor{blue}{14,3 $\pm$ 2,2 pp} \\
        \textcolor{blue}{SimpleRNN} & \textcolor{blue}{66,4\% $\pm$ 1,8 pp} & \textcolor{blue}{9,8 $\pm$ 3,0 pp} \\
        \textcolor{blue}{GRU} & \textcolor{blue}{66,2\% $\pm$ 0,6 pp} & \textcolor{blue}{10,7 $\pm$ 0,5 pp} \\
        \textcolor{blue}{CNN+LSTM} & \textcolor{blue}{65,7\% $\pm$ 1,3 pp} & \textcolor{blue}{16,2 $\pm$ 0,9 pp} \\
        \bottomrule
    \end{tabular}
\end{table}

\begin{figure}[H]
    \centering
    \includegraphics[width=\textwidth]{rnn_loss_compare}
    \caption{\textcolor{blue}{Acurácia de validação e \textit{gap} treino/validação ao final do treinamento, média $\pm$ desvio padrão sobre 5 sementes independentes por arquitetura, sob divisão temporal.}}
    \label{fig:res_rnn_loss}
\end{figure}

\subsection{Robustez: validação cruzada, matrizes de confusão e separabilidade das classes}
\label{sec:resultados_robustez}

Esta seção final de resultados responde a três perguntas que as seções anteriores deixam em aberto: onde os classificadores erram, não apenas o quanto acertam; se as acurácias reportadas são estáveis sob reamostragem; e como as classes se distribuem visualmente no espaço de atributos.

\textcolor{blue}{A Tabela~\ref{tab:res_cv} compara os seis modelos clássicos sob validação cruzada temporal em blocos de 5 dobras (Seção~\ref{sec:protocolo}), para os dois problemas -- 17 classes e o subconjunto de 6 classes mecânicas puras (Seção~\ref{sec:resultados_sensibilidade}). Random Forest lidera nas 17 classes (F1-macro 0,913); SVM lidera no subconjunto mecânico (F1-macro 0,798). Diferente da versão preliminar deste protocolo -- que reportava dispersão entre dobras inferior a 0,6 ponto percentual sob \texttt{StratifiedKFold} aleatório --, a dispersão sob dobras temporais em blocos é substancialmente maior: entre 3,2 e 4,3 pontos percentuais de F1-macro nas 17 classes, e entre 5,4 e 9,1 pontos percentuais nas 6 classes mecânicas. Isso não é um defeito do novo protocolo; é o preço, em variância, de dobras que não são mais permutáveis entre si por construção -- cada dobra testa um bloco contíguo e distinto da sessão de cada rótulo, então dobras diferentes podem genuinamente corresponder a trechos da sessão com comportamento espectral algo diferente, o que uma reamostragem aleatória (ao misturar aquisições de toda a sessão em cada dobra) mascarava.}

\begin{table}[H]
    \centering
    \caption{\textcolor{blue}{Acurácia e F1-macro sob validação cruzada temporal em blocos de 5 dobras (Seção~\ref{sec:protocolo}).}}
    \label{tab:res_cv}
    \begin{tabular}{@{}llrr@{}}
        \toprule
        \textbf{Problema} & \textbf{Modelo} & \textbf{Acurácia (5-fold)} & \textbf{F1-macro (5-fold)} \\
        \midrule
        \textcolor{blue}{17 classes} & \textcolor{blue}{Random Forest} & \textcolor{blue}{89,10\% $\pm$ 5,46\%} & \textcolor{blue}{0,913 $\pm$ 0,043}\\
        \textcolor{blue}{17 classes} & \textcolor{blue}{k-NN} & \textcolor{blue}{88,34\% $\pm$ 4,32\%} & \textcolor{blue}{0,906 $\pm$ 0,038}\\
        \textcolor{blue}{17 classes} & \textcolor{blue}{SVM} & \textcolor{blue}{87,97\% $\pm$ 3,89\%} & \textcolor{blue}{0,902 $\pm$ 0,035}\\
        \textcolor{blue}{17 classes} & \textcolor{blue}{MLP} & \textcolor{blue}{89,51\% $\pm$ 3,62\%} & \textcolor{blue}{0,899 $\pm$ 0,032}\\
        \textcolor{blue}{17 classes} & \textcolor{blue}{Regressão Logística} & \textcolor{blue}{89,03\% $\pm$ 3,91\%} & \textcolor{blue}{0,898 $\pm$ 0,040}\\
        \textcolor{blue}{17 classes} & \textcolor{blue}{XGBoost} & \textcolor{blue}{89,72\% $\pm$ 4,28\%} & \textcolor{blue}{0,898 $\pm$ 0,037}\\
        \textcolor{blue}{6 classes} & \textcolor{blue}{SVM} & \textcolor{blue}{80,11\% $\pm$ 5,61\%} & \textcolor{blue}{0,798 $\pm$ 0,058}\\
        \textcolor{blue}{6 classes} & \textcolor{blue}{Random Forest} & \textcolor{blue}{78,80\% $\pm$ 5,97\%} & \textcolor{blue}{0,784 $\pm$ 0,063}\\
        \textcolor{blue}{6 classes} & \textcolor{blue}{k-NN} & \textcolor{blue}{78,35\% $\pm$ 5,29\%} & \textcolor{blue}{0,781 $\pm$ 0,054}\\
        \textcolor{blue}{6 classes} & \textcolor{blue}{MLP} & \textcolor{blue}{79,68\% $\pm$ 6,87\%} & \textcolor{blue}{0,764 $\pm$ 0,069}\\
        \textcolor{blue}{6 classes} & \textcolor{blue}{Regressão Logística} & \textcolor{blue}{78,67\% $\pm$ 7,59\%} & \textcolor{blue}{0,759 $\pm$ 0,091}\\
        \textcolor{blue}{6 classes} & \textcolor{blue}{XGBoost} & \textcolor{blue}{77,83\% $\pm$ 7,64\%} & \textcolor{blue}{0,720 $\pm$ 0,078}\\
        \bottomrule
    \end{tabular}
\end{table}

\textcolor{blue}{A Figura~\ref{fig:res_cv_scatter} exibe a acurácia individual de cada uma das cinco dobras temporais em blocos (pontos dispersos, com leve deslocamento horizontal para legibilidade) junto da média com barra de erro (marcador contornado em preto), para SVM e Regressão Logística, nos dois problemas (17 classes e 6 classes mecânicas). A dispersão entre dobras sob o protocolo corrigido é substancialmente maior do que a sugerida pela versão preliminar desta figura, sobretudo no subconjunto mecânico: a dobra 0 do SVM (70,1\% de acurácia) e a dobra 4 da Regressão Logística (65,2\%) destoam visivelmente das demais dobras do mesmo modelo (82--83\%), consistente com a explicação já registrada no texto -- dobras temporais em blocos testam trechos genuinamente distintos da sessão de cada rótulo, ao contrário de uma reamostragem aleatória, que mistura aquisições de toda a sessão em cada dobra e mascara essa variação.}

\begin{figure}[H]
    \centering
    \includegraphics[width=0.75\textwidth]{addendum_cv}
    \caption{Acurácia por dobra individual (pontos dispersos) e média $\pm$ desvio padrão (marcador com barra de erro), sob divisão temporal em blocos, para SVM e Regressão Logística nos dois problemas (17 e 6 classes).}
    \label{fig:res_cv_scatter}
\end{figure}

A matriz de confusão do classificador SVM para as 17 classes (Figura~\ref{fig:res_cm17}, normalizada por linha) mostra diagonal fortemente dominante para a quase totalidade das condições. A única confusão visualmente relevante ocorre entre \texttt{MÁQUINA BALANCEADA (S/~PARAF)} e \texttt{MÁQUINA DESBALANCEADA (RET~14~PARAF)}. Restringindo a mesma análise às seis classes mecânicas puras (Figura~\ref{fig:res_cm6}), essa confusão é quantificada diretamente: 61\% das amostras de \texttt{MÁQUINA BALANCEADA} são classificadas como \texttt{RET~14~PARAF}, e 53\% das amostras de \texttt{RET~14~PARAF} são classificadas como \texttt{MÁQUINA BALANCEADA} -- muito acima de qualquer outro par de classes, cujas taxas de acerto individuais chegam a 99--100\%. Como registrado na Seção~\ref{sec:resultados_bd}, o rótulo \texttt{RET~N~PARAF} é literalmente a quantidade de parafusos retirados de um disco calibrado para graduar o desbalanceamento -- \texttt{S/~PARAF} e \texttt{RET~14~PARAF} são, portanto, dois graus quantitativos do mesmo tipo de falha, não duas condições qualitativamente distintas, o que confere plausibilidade física direta à proximidade observada na assinatura espectral e explica por que essa confusão persiste apesar dos elevados tamanhos de efeito ($\eta^2$) reportados na Seção~\ref{sec:resultados_sensibilidade} para o conjunto das seis classes mecânicas como um todo.

\begin{figure}[H]
    \centering
    \includegraphics[width=0.85\textwidth]{addendum_cm17}
    \caption{Matriz de confusão, classificador SVM, 17 classes (normalizada por linha).}
    \label{fig:res_cm17}
\end{figure}

\begin{figure}[H]
    \centering
    \includegraphics[width=0.62\textwidth]{addendum_cm6}
    \caption{Matriz de confusão, classificador SVM, 6 classes mecânicas puras (normalizada por linha, valores anotados em \%).}
    \label{fig:res_cm6}
\end{figure}

Por fim, a projeção em duas componentes principais do espaço de atributos selecionado (Figura~\ref{fig:res_pca_full}, 1.000 amostras, 40 atributos, coloridas pelas três categorias da Tabela~\ref{tab:res_taxonomia}) mostra separação visual quase completa entre condições mecânicas puras, elétricas puras e combinadas: a primeira componente principal (PC1) explica 57,9\% da variância total, e já é suficiente para isolar o agrupamento mecânico (à esquerda) dos agrupamentos elétricos (à direita, subdivididos em múltiplos sub-agrupamentos correspondentes às diferentes condições \texttt{DIODO}); a segunda componente (PC2, 8,7\%) separa o agrupamento combinado do agrupamento mecânico puro. Restringindo a mesma projeção às seis condições mecânicas (Figura~\ref{fig:res_pca_mech}, PC1 = 60,7\%, PC2 = 12,1\%), a condição \texttt{MÁQUINA DESALINHADA} forma um agrupamento nitidamente isolado, enquanto as demais cinco condições -- entre elas o par \texttt{S/~PARAF}/\texttt{RET~14~PARAF} já discutido -- apresentam sobreposição parcial visualmente consistente com a acurácia de 80--81\% da Tabela~\ref{tab:res_cv}: separável, mas não trivialmente.

\begin{figure}[H]
    \centering
    \includegraphics[width=0.62\textwidth]{addendum_pca_full}
    \caption{Projeção em duas componentes principais (1.000 amostras, 40 atributos selecionados), coloridas por categoria da taxonomia (Tabela~\ref{tab:res_taxonomia}).}
    \label{fig:res_pca_full}
\end{figure}

\begin{figure}[H]
    \centering
    \includegraphics[width=0.62\textwidth]{addendum_pca_mech}
    \caption{Projeção em duas componentes principais restrita às 6 condições mecânicas puras (1.000 amostras).}
    \label{fig:res_pca_mech}
\end{figure}

\subsection{Síntese dos achados frente às perguntas de pesquisa}
\label{sec:resultados_sintese}

\textcolor{blue}{Em relação a RQ1, o conjunto de evidências reunidas neste capítulo -- acurácia de 88--90\% sob validação cruzada temporal nas 17 condições (Tabela~\ref{tab:res_cv}), convergência entre importância SHAP e tamanho de efeito estatístico ($\eta^2$) na sub-análise de sensibilidade mecânica (Seção~\ref{sec:resultados_sensibilidade}), a ablação controlada H/IH/H+IH que confirma diretamente a contribuição das inter-harmônicas para o subconjunto mecânico (Seção~\ref{sec:resultados_ablacao}), e separação visual das três categorias de falha no espaço de componentes principais (Figura~\ref{fig:res_pca_full}) -- sustenta que um modelo treinado sobre a assinatura espectral agregada diagnostica corretamente o modo de falha do sistema trifásico já instalado, com desempenho suficiente para apoiar decisões de manutenção baseada em condição; não foi testada, neste estágio, a antecedência com que um modo de falha em desenvolvimento seria detectado (ver o baseline de persistência discutido logo abaixo, junto à Figura~\ref{fig:res_forecast}), tal como formulado na Seção~\ref{sec:objetivos} do Capítulo~\ref{cap:fundamentacao}.} Em relação a RQ2, a evidência corrigida da Seção~\ref{sec:resultados_rq2} (Figura~\ref{fig:res_nilm}) sustenta viabilidade real, porém parcial e seletiva por canal, da desagregação de cargas a partir da mesma assinatura espectral -- com o achado adicional, não antecipado pela literatura de NILM revisada no Capítulo~\ref{cap:fundamentacao}, de que canais inter-harmônicos se mostraram mais recuperáveis que harmônicos entre os dez melhores canais avaliados. A interpretação crítica conjunta desses achados frente à literatura revisada, suas implicações para as duas perguntas de pesquisa e suas limitações estruturais são desenvolvidas no Capítulo~\ref{cap:discussao} e no Capítulo~\ref{cap:limitacoes}, respectivamente.
%===============================================================
\section{Discussão}
\label{cap:discussao}

Este capítulo interpreta criticamente os resultados apresentados no Capítulo~\ref{cap:resultados} à luz da literatura revisada no Capítulo~\ref{cap:fundamentacao}, articulando-os diretamente às duas perguntas de pesquisa formuladas na Seção~\ref{sec:objetivos}. A Seção~\ref{sec:discussao_rq1} discute a evidência de RQ1, integrando o desempenho de classificação, a sub-análise de sensibilidade mecânica isolada e a interpretabilidade via SHAP. A Seção~\ref{sec:discussao_rq2} discute a evidência de RQ2, com ênfase no papel inesperado das componentes interharmônicas na desagregação de cargas. A Seção~\ref{sec:discussao_contribuicoes} sintetiza as contribuições metodológicas da pesquisa frente às lacunas identificadas no Capítulo~\ref{cap:fundamentacao}. A Seção~\ref{cap:limitacoes} discute as limitações estruturais do estudo, condição necessária para delimitar corretamente o alcance das conclusões apresentadas no Capítulo~\ref{cap:conclusao}.

\subsection{RQ1 -- diagnóstico de falha: convergência entre desempenho, interpretabilidade e literatura de MCSA}
\label{sec:discussao_rq1}

\textcolor{blue}{O conjunto de evidências reunido na Seção~\ref{sec:resultados_rq1} responde afirmativamente a RQ1 dentro do escopo testado: a acurácia de 87,97\%--89,72\% (a depender do modelo) sob validação cruzada temporal em blocos nas 17 condições (Tabela~\ref{tab:res_cv}), com dispersão entre dobras de 3 a 4 pontos percentuais de F1-macro (Figura~\ref{fig:res_cv_scatter} -- maior do que a dispersão sob reamostragem aleatória, pelas razões discutidas na Seção~\ref{sec:protocolo}), indica que o modelo treinado sobre a assinatura espectral agregada discrimina corretamente o modo de falha também quando testado contra o trecho final, cronologicamente posterior, de cada sessão, não apenas contra aquisições vizinhas às de treino. A classificação hierárquica das quatro macro-categorias (Seção~\ref{sec:resultados_macro}) qualifica essa resposta de forma direta: quando o problema é reduzido a distinguir apenas saudável, elétrica isolada, mecânica isolada e combinada, o desempenho sobe para até 97,97\% de acurácia (Random Forest, Tabela~\ref{tab:res_macro}) -- evidência de que a dificuldade do problema de 17 classes está concentrada na distinção de subtipo e severidade dentro de cada ramo, não na distinção entre os quatro macroestados em si, que a assinatura espectral separa quase perfeitamente. Essa mesma reformulação, porém, expõe uma dispersão entre dobras temporais muito maior (14--18 pontos percentuais de F1-macro) do que a das 17 classes, atribuível ao reduzido número de amostras da categoria \texttt{saudável} (uma única condição, 1.440 amostras) diluído em blocos temporais estreitos -- uma qualificação que deve acompanhar a leitura otimista dos quase 98\% de acurácia macro.} Esse resultado corrobora empiricamente, para um espaço de atributos harmônico e interharmônico bem mais amplo, a premissa central da MCSA clássica -- de que a corrente de estator, já disponível sem instrumentação adicional, carrega informação diagnóstica suficiente para sustentar decisões de manutenção \citep{thomson2001,nandi2005}. A importância SHAP do classificador treinado sobre as 17 classes (Figura~\ref{fig:res_shap_full}), dominada por harmônicos de corrente com a tensão aparecendo apenas na décima posição, reforça essa leitura: o modelo, sem qualquer restrição \textit{a priori} de domínio, redescobre por conta própria a mesma preferência por corrente que a MCSA clássica assume por construção.

\textcolor{blue}{Essa concordância entre modelo e teoria, contudo, esconde uma heterogeneidade importante que só emerge na sub-análise de sensibilidade mecânica isolada (Seção~\ref{sec:resultados_sensibilidade}). Quando as seis condições de falha elétrica são removidas do problema e apenas as seis condições mecânicas puras são avaliadas, o F1-macro do SVM cai cerca de 7 pontos percentuais (de 0,898 para 0,827, Tabelas~\ref{tab:res_classical} e~\ref{tab:res_classical_mech})} e a importância SHAP se inverte: o atributo mais relevante passa a ser de tensão, e cinco dos quinze atributos mais importantes são de tensão, contra apenas um de quinze no problema completo. \textcolor{blue}{Essa inversão sugere maior conteúdo discriminativo em atributos de tensão para falha mecânica do que para falha elétrica; uma interpretação fisicamente plausível é que desbalanceamento e desalinhamento afetam primariamente o entreferro e a força eletromotriz induzida, mas essa leitura permanece hipótese, não fato demonstrado por esta pesquisa -- a tensão no ponto de acoplamento comum também depende da impedância da fonte, da rede e da instrumentação, e o mecanismo eletromecânico exato precisa de investigação dedicada, considerando a topologia e a impedância da alimentação}, e estende empiricamente, para um espaço de atributos espectrais amplo e não restrito a bandas laterais específicas, uma distinção que a literatura de MCSA revisada trata de forma mais implícita: \citet{nandi2005} sistematizam os mecanismos físicos de diferentes modos de falha sem quantificar a contribuição relativa de corrente e tensão para cada um deles, e nem \citet{benamira2022} nem \citet{martinez2022} comparam explicitamente a relevância de tensão \textit{versus} corrente entre falha elétrica e falha mecânica isolada dentro de um mesmo protocolo experimental. A convergência entre SHAP (baseado em contribuição marginal ao modelo) e $\eta^2$ (baseado em tamanho de efeito estatístico, calculado de forma totalmente independente do classificador) para essa mesma inversão tensão/corrente é, nesse sentido, uma evidência metodologicamente mais robusta do que qualquer um dos dois indicadores isoladamente forneceria.

O achado mais diretamente conectado à motivação teórica desta pesquisa, contudo, está no lado interharmônico da sub-análise de sensibilidade: as componentes de ordem 0 e 1 da corrente (aproximadamente contínuas) concentram os maiores tamanhos de efeito de todo o levantamento ($\eta^2$ entre 0,90 e 0,95), superando o maior $\eta^2$ harmônico observado (0,863). Esse resultado estende diretamente, para o problema de sensibilidade mecânica, a motivação estabelecida por \citet{efe2016} de que componentes interharmônicas podem carregar distorção percentual superior à da harmônica adjacente -- motivação originalmente formulada para um sistema retificador puramente elétrico, e aqui confirmada, por uma via empírica distinta, também para falha mecânica isolada. Trata-se de uma extensão genuína da lacuna identificada na Seção~\ref{sec:sintese} do Capítulo~\ref{cap:fundamentacao}: nenhum dos doze trabalhos de NILM e de MCSA ali revisados avalia sistematicamente o conteúdo interharmônico como atributo diagnóstico, e o presente resultado indica que essa omissão pode ter descartado, também para diagnóstico mecânico, um dos sinais espectrais mais informativos disponíveis.

A matriz de confusão do classificador SVM restrita às seis classes mecânicas (Figura~\ref{fig:res_cm6}) qualifica essa evidência de forma relevante para a interpretação de RQ1 como problema de classificação: a confusão dominante -- 61\% de \texttt{MÁQUINA BALANCEADA} classificada como \texttt{RET~14~PARAF} e 53\% no sentido inverso -- não é um erro aleatório do classificador, mas reflexo direto de uma característica do desenho experimental, já registrada na Seção~\ref{sec:resultados_bd}: o rótulo \texttt{RET~N~PARAF} corresponde ao número de parafusos retirados de um disco calibrado, ou seja, é uma escala quantitativa de grau de desbalanceamento, não uma taxonomia de falhas qualitativamente distintas. Esse ponto tem implicação metodológica direta para trabalhos futuros: \citet{martinez2022} e \citet{benamira2022} tratam a detecção de desbalanceamento mecânico como um problema binário (presente \textit{versus} ausente), e o presente resultado sugere que, para condições de desbalanceamento de baixa severidade, uma formulação de classificação discreta tende a subestimar a real capacidade diagnóstica do sistema -- uma reformulação ordinal ou de regressão sobre o grau de desbalanceamento, e não apenas sobre sua presença, poderia capturar de forma mais fiel a informação já presente na assinatura espectral, hipótese que esta pesquisa não testou diretamente e que é retomada na Seção~\ref{cap:limitacoes}.

\textcolor{blue}{Por fim, a comparação entre as seis arquiteturas sequenciais (Seção~\ref{sec:resultados_rnn}) endereça diretamente a lacuna apontada por \citet{hoang2019} quanto à escassez de redes neurais recorrentes aplicadas ao diagnóstico de falha de máquinas rotativas, estendendo-a para incluir também uma arquitetura baseada em atenção \citep{vaswani2017}, fora do escopo da revisão original de \citeauthor{hoang2019}. Sob o protocolo corrigido -- divisão temporal 75/25 por rótulo aplicada antes do janelamento (Seção~\ref{sec:protocolo}), eliminando o compartilhamento de aquisições entre janelas de treino e de validação apontado como vazamento pelo orientador --, uma primeira execução de semente única indicou o Transformer como a arquitetura de pior desempenho; a repetição sob cinco sementes independentes (Tabela~\ref{tab:res_rnn_multiseed}) mostrou que essa conclusão não se sustentava -- o Transformer é, na verdade, a arquitetura de maior acurácia média de validação (68,6\% $\pm$ 2,0 pontos percentuais), seguido por LSTM (67,2\%), BiLSTM (66,8\%), SimpleRNN (66,4\%), GRU (66,2\%) e CNN+LSTM (65,7\%), com desvios-padrão entre sementes (0,6--3,0 pontos percentuais) grandes o suficiente para que a maioria dessas diferenças não deva ser lida como estabelecida além do acaso. Esse episódio é, em si, o argumento mais direto reunido nesta pesquisa a favor de comparações multi-semente para arquiteturas de aprendizado profundo: uma execução isolada teria levado a uma conclusão qualitativamente oposta -- e incorreta -- sobre qual arquitetura generaliza melhor. Um padrão mais estável emergiu, porém, do próprio \textit{gap} treino/validação: o Transformer manteve o maior \textit{gap} médio nas cinco sementes (20,6 $\pm$ 1,1 pontos percentuais, contra 9,8--16,2 pontos percentuais nas demais), ou seja, é consistentemente a arquitetura que mais se ajusta aos dados de treino -- mas isso não a impede de generalizar, em média, tão bem ou melhor do que as demais em termos absolutos. A leitura mais honesta desse padrão é que a sofisticação do mecanismo de atenção não determina, isoladamente, nem vantagem nem desvantagem de generalização neste espaço de atributos: o Transformer sobre-ajusta mais e, ainda assim, generaliza pelo menos tão bem quanto as arquiteturas recorrentes e convolucional, um resultado que só a repetição multi-semente (Seção~\ref{cap:limitacoes}) tornou possível distinguir de ruído de inicialização. Uma ressalva crítica, porém, precisa ser registrada: as arquiteturas sequenciais foram comparadas sob um protocolo distinto do usado para os modelos clássicos (janela de 8 aquisições sobre 20 atributos, contra 40 atributos selecionados e divisão temporal completa para SVM, Regressão Logística e demais modelos clássicos da Tabela~\ref{tab:res_classical}), de modo que a acurácia média máxima obtida pelas RNNs sob o protocolo corrigido (68,6\%, Transformer) fica bem abaixo da acurácia dos modelos clássicos sob o mesmo tipo de divisão temporal (87--90\%, Tabela~\ref{tab:res_classical}) -- uma diferença que não deve ser lida como evidência de que modelos clássicos superam arquiteturas sequenciais para este problema em geral, mas como resultado direto de um espaço de atributos vinte vezes menor (20 \textit{versus} 40 atributos) e um número de janelas de treino reduzido (3.635) para as RNNs, uma limitação de protocolo discutida na Seção~\ref{cap:limitacoes}. A comparação entre as seis arquiteturas sequenciais entre si, obtida sob protocolo idêntico, permanece, isso sim, diretamente interpretável -- e é precisamente essa comparação interna que só se estabilizou sob repetição multi-semente.}

\subsection{RQ2 -- desagregação de cargas: viabilidade parcial e o papel inesperado das interharmônicas}
\label{sec:discussao_rq2}

\textcolor{blue}{A evidência corrigida de RQ2 (Seção~\ref{sec:resultados_rq2}, Figura~\ref{fig:res_nilm}), agora com seleção de atributos sem vazamento (Seção~\ref{sec:protocolo}), sustenta uma resposta afirmativa, porém qualificada: é possível desagregar, a partir exclusivamente da assinatura espectral agregada e sem forma de onda bruta, contribuições espectrais individuais de cargas operando simultaneamente -- mas apenas para um subconjunto seletivo de canais, com o melhor canal individual atingindo $R^2=0{,}158$ e 86,7\% dos 60 canais avaliados apresentando $R^2$ negativo. A ablação controlada (Seção~\ref{sec:resultados_ablacao}) qualifica essa evidência de forma ainda mais direta: usando somente harmônicas, nenhum canal atinge $R^2$ positivo; usando somente inter-harmônicas, um canal atinge $R^2=0{,}880$ -- o maior valor de toda a análise de RQ2.} Essa limitação de desempenho é consistente com o mecanismo identificado por \citet{bouhouras2017}: como a desagregação nesta pesquisa, assim como na de \citeauthor{bouhouras2017}, opera exclusivamente sobre amplitudes espectrais agregadas, sem acesso à fase relativa entre as componentes de diferentes cargas, a soma vetorial de correntes harmônicas de múltiplas fontes com conteúdo espectral significativo não se reduz a uma soma escalar simples -- limitação estrutural do próprio paradigma de desagregação baseado em amplitude agregada, não uma falha específica do modelo empregado nesta pesquisa. Sob essa leitura, o desempenho seletivo por canal observado aqui é, na verdade, o resultado esperado desse mecanismo, e não uma anomalia a ser corrigida por ajuste fino de hiperparâmetros.

\textcolor{blue}{O achado mais relevante para a segunda metade de RQ2 -- a contribuição específica das componentes interharmônicas -- é o predomínio de canais interharmônicos entre os mais recuperáveis: dois dos quatro melhores canais são interharmônicos ($IH_{46}\_I_c\_avg$ e $IH_{48}\_I_c\_avg$, os dois primeiros colocados; os outros dois são $THDe\_I_c\_avg$ e $H_{12}\_I_b\_avg$) e catorze dos vinte melhores são interharmônicos, apesar de a etapa de seleção de atributos sem vazamento (Seção~\ref{sec:protocolo}) ter produzido uma distribuição próxima do equilíbrio entre harmônicas e interharmônicas (28 e 29, entre os 60 canais selecionados), o que descarta a hipótese de que esse predomínio seja artefato de um viés na pré-seleção.} Esse resultado é diretamente contraintuitivo frente aos três trabalhos de NILM revisados no Capítulo~\ref{cap:fundamentacao} -- \citet{djordjevic2017}, \citet{bouhouras2017} e \citet{li2019moea} restringem suas assinaturas de carga a ordens harmônicas inteiras, sem testar explicitamente componentes interharmônicas --, e constitui a confirmação mais direta, entre todos os resultados desta pesquisa, da hipótese originalmente levantada por \citet{efe2016} de que a exclusão sistemática de interharmônicas nas formulações de assinatura de carga pode representar descarte de informação espectral discriminativa. É também, nesse sentido, o achado que mais diretamente materializa a segunda pergunta de pesquisa formulada na Seção~\ref{sec:objetivos}: não apenas as interharmônicas contribuem para a desagregação, como se mostram, entre os canais avaliados, comparativamente mais informativas do que as harmônicas correspondentes.

Essa evidência, no entanto, precisa ser lida à luz de duas ressalvas registradas na Seção~\ref{sec:resultados_rq2}. Primeira, a correção metodológica do próprio cálculo de $R^2$: sob validação \textit{leave-one-condition-out} com apenas quatro condições combinadas disponíveis, o alvo é constante dentro de cada dobra, de modo que o $R^2$ por dobra individual degenera para o caso especial de variância nula -- um resultado que, antes de ser corrigido, poderia ter sido mal interpretado como desempenho nulo ou perfeito sem qualquer relação com a real capacidade de generalização do modelo. Essa correção, embora não seja em si um resultado substantivo de RQ2, constitui uma contribuição metodológica relevante para qualquer estudo futuro de desagregação de cargas que, como este, opere sob um número reduzido de condições combinadas e valide por retenção de condição completa: a métrica de generalização real exige agrupar as previsões de todas as dobras antes de pontuar, não pontuar dobra a dobra. Segunda, o tamanho amostral de apenas quatro condições combinadas -- consequência direta do desenho experimental da bancada de retificador (Seção~\ref{sec:bancadas}), não de uma limitação do pipeline de análise -- levou à opção por um \textit{Random Forest Regressor} em vez de uma arquitetura sequencial para RQ2, em contraste direto com as seis arquiteturas comparadas para RQ1. Essa assimetria de profundidade analítica entre as duas perguntas de pesquisa é discutida como limitação estrutural na Seção~\ref{cap:limitacoes}, e delimita RQ2, no estágio atual desta pesquisa, como um estudo de viabilidade, não como um sistema de desagregação pronto para aplicação irrestrita.

\subsection{Contribuições metodológicas frente às lacunas identificadas na literatura}
\label{sec:discussao_contribuicoes}

Para além das respostas específicas a RQ1 e RQ2, três contribuições metodológicas desta pesquisa merecem registro explícito por endereçarem diretamente lacunas identificadas na síntese crítica da literatura (Seção~\ref{sec:sintese} do Capítulo~\ref{cap:fundamentacao}). Primeira, nenhum dos doze trabalhos de NILM e de MCSA ali revisados trata diagnóstico de falha mecânica e desagregação de cargas elétricas como duas saídas de um mesmo espaço de entrada espectral agregado em um sistema trifásico industrial; esta pesquisa demonstra, ainda que com desempenho desigual entre as duas tarefas, que uma mesma assinatura espectral -- harmônica e interharmônica, de tensão e corrente, extraída do mesmo ponto de acoplamento comum -- pode alimentar simultaneamente um classificador de 17 classes (RQ1) e um regressor multi-saída de desagregação (RQ2), sem instrumentação ou pipeline de aquisição distintos entre os dois problemas. Segunda, a relevância das componentes interharmônicas, motivada isoladamente por \citet{efe2016} em um único cenário retificador, foi confirmada empiricamente em três contextos analíticos distintos e independentes desta pesquisa -- a sub-análise de sensibilidade mecânica ($\eta^2$ de até 0,95 nas ordens 0 e 1), a desagregação de cargas ($R^2$ predominantemente maior entre canais interharmônicos que harmônicos entre os dez melhores) e, de forma indireta, a própria decisão metodológica de ampliar o espaço de atributos além do conjunto inicialmente restrito a correntes (Seção~\ref{sec:bd}) -- reforçando com evidência convergente, e não apenas replicada, uma hipótese que a literatura de NILM revisada nunca havia testado sistematicamente. Terceira, a convergência entre dois métodos de avaliação de importância de atributos metodologicamente independentes -- SHAP \citep{lundberg2017}, baseado em contribuição marginal a um modelo treinado, e $\eta^2$/delta de Cliff, baseados em tamanho de efeito estatístico calculado sem qualquer modelo preditivo -- para o mesmo padrão de inversão tensão/corrente entre falha elétrica e falha mecânica fortalece a confiança na interpretação física dos resultados além do que qualquer um dos dois métodos, isoladamente, permitiria afirmar.

\subsection{Limitações}
\label{cap:limitacoes}

\textcolor{blue}{Cinco} limitações estruturais delimitam o alcance das conclusões desta pesquisa e devem ser explicitamente registradas antes de sua generalização. Primeira, defeitos em rolamentos -- previstos no desenho experimental e amplamente documentados como um dos modos de falha mais estudados na literatura de diagnóstico de máquinas rotativas \citep{hoang2019} -- ainda não foram incorporados à base de dados analisada; trata-se de uma lacuna de coleta, não de processamento, e sua ausência significa que a taxonomia de 17 condições, embora ampla, não cobre a totalidade dos modos de falha mecânica tipicamente considerados nesse campo. Segunda, a evidência de RQ2 baseia-se em apenas quatro condições combinadas, o que não apenas caracteriza a desagregação de cargas como um estudo de viabilidade (Seção~\ref{sec:discussao_rq2}) como também impediu o emprego de arquiteturas sequenciais comparáveis às utilizadas em RQ1, introduzindo uma assimetria de profundidade analítica entre as duas perguntas de pesquisa que trabalhos futuros, com bancadas capazes de gerar mais combinações de carga, poderiam corrigir. Terceira, a escala de severidade do desbalanceamento mecânico (\texttt{RET~N~PARAF}) não foi fisicamente caracterizada nesta pesquisa -- massa e geometria dos parafusos retirados do disco calibrado não foram quantificadas em termos de momento de desbalanceamento --, o que limita a interpretação quantitativa (embora não a qualitativa) da confusão entre \texttt{MÁQUINA BALANCEADA} e \texttt{RET~14~PARAF} discutida na Seção~\ref{sec:discussao_rq1} e restringe a uma hipótese, ainda não testada, a reformulação ordinal desse problema sugerida naquela seção. \textcolor{blue}{Quarta -- parcialmente corrigida ao longo desta pesquisa: as curvas de treinamento das seis arquiteturas sequenciais (Figuras~\ref{fig:res_rnn_curves}--\ref{fig:res_rnn_loss}) ainda ilustram apenas a execução de semente única originalmente reportada, mas a comparação de desempenho em si (Tabela~\ref{tab:res_rnn_multiseed}, Seção~\ref{sec:resultados_rnn}) foi refeita sob cinco sementes independentes por arquitetura, e o resultado justificou plenamente a cautela desta limitação: a ordenação de semente única (Transformer como pior arquitetura) inverteu-se sob repetição (Transformer como melhor arquitetura em média). Ainda assim, cinco sementes é uma amostra modesta frente às cinco dobras de validação cruzada temporal disponíveis para os modelos clássicos (Tabela~\ref{tab:res_cv}), de modo que os desvios-padrão entre sementes reportados para as RNNs devem ser lidos como estimativas menos precisas de variância do que os das dobras clássicas, e diferenças pequenas entre arquiteturas adjacentes (por exemplo, entre LSTM e BiLSTM) continuam exigindo cautela interpretativa adicional.} \textcolor{blue}{Quinta, e possivelmente a mais relevante das cinco para a generalização dos resultados}, a totalidade dos resultados desta pesquisa foi obtida em duas bancadas experimentais específicas, sem replicação independente em equipamentos ou instalações distintas -- uma limitação comum a estudos de diagnóstico e desagregação de escopo experimental controlado, mas que não deve ser omitida ao se avaliar a generalização dos resultados para outros sistemas trifásicos industriais.

\section{Conclusão}
\label{cap:conclusao}

Esta pesquisa investigou se uma mesma assinatura espectral agregada -- componentes harmônicas e interharmônicas de tensão e corrente, obtidas de analisadores de energia já presentes em instalações industriais, sem forma de onda bruta nem instrumentação adicional dedicada -- é suficiente para responder simultaneamente a duas perguntas de pesquisa co-primárias: o diagnóstico de modo de falha em sistemas trifásicos (RQ1) e a desagregação de cargas elétricas operando simultaneamente no mesmo ponto de acoplamento comum (RQ2).

Em relação a RQ1, o conjunto de evidências reunido no Capítulo~\ref{cap:resultados} e discutido na Seção~\ref{sec:discussao_rq1} sustenta uma resposta afirmativa: um modelo treinado sobre a assinatura espectral agregada diagnostica corretamente o modo de falha do sistema trifásico -- saudável, falha elétrica isolada, falha mecânica isolada ou falha combinada, entre 17 condições experimentais -- \textcolor{blue}{com acurácia de 88--90\% sob validação cruzada temporal, desempenho suficiente para diagnosticar o modo de falha já instalado e apoiar decisões de manutenção baseada em condição -- não, neste estágio, para detecção precoce de uma falha em desenvolvimento, cuja evidência preliminar (previsão de um passo à frente) não superou um preditor de persistência trivial (Seção~\ref{sec:resultados_rq1}).} Essa evidência é reforçada por \textcolor{blue}{cinco} linhas convergentes e metodologicamente independentes: \textcolor{blue}{a classificação hierárquica das quatro macro-categorias (Seção~\ref{sec:resultados_macro}), que isola a dificuldade real do problema de 17 classes na distinção de subtipo e severidade dentro de cada ramo -- não entre os ramos em si, separados com até 97,97\% de acurácia (Tabela~\ref{tab:res_macro});} a interpretabilidade via SHAP, que redescobre sem restrição \textit{a priori} a preferência por corrente já assumida pela MCSA clássica \citep{thomson2001,nandi2005}; a sub-análise de sensibilidade mecânica isolada, que revela tanto a maior sutileza da assinatura mecânica \textcolor{blue}{(queda de 7 pontos percentuais de F1-macro ao isolar as seis condições mecânicas)} quanto a inversão física tensão/corrente entre falha elétrica e mecânica; \textcolor{blue}{a ablação controlada H/IH/H+IH (Seção~\ref{sec:resultados_ablacao}), que confirma que as inter-harmônicas sozinhas superam as harmônicas sozinhas por quase 11 pontos percentuais de F1-macro no subconjunto mecânico;} \textcolor{blue}{e a comparação sistemática de seis arquiteturas sequenciais de aprendizado profundo sob divisão temporal e repetição multi-semente (Seção~\ref{sec:resultados_rnn}), que endereça diretamente a escassez de redes neurais recorrentes aplicadas a esse domínio apontada por \citet{hoang2019}: as seis arquiteturas generalizaram de forma relativamente próxima entre si (65,7--68,6\% de acurácia média de validação, com desvios-padrão entre sementes que se sobrepõem para a maioria dos pares), com o Transformer temporal \citep{vaswani2017} à frente em média apesar de apresentar também o maior \textit{gap} treino/validação das seis de forma consistente -- um resultado que reverteu, sob repetição multi-semente, a ordenação oposta observada na execução inicial de semente única, e que ilustra diretamente por que a comparação entre arquiteturas profundas exige múltiplas sementes para ser interpretável}.

\textcolor{blue}{Em relação a RQ2, a evidência corrigida discutida na Seção~\ref{sec:discussao_rq2} sustenta viabilidade real, porém parcial e seletiva por canal: o melhor canal espectral avaliado atinge $R^2=0{,}158$, e a desagregação não deve ser lida, no estágio atual desta pesquisa, como um sistema pronto para aplicação irrestrita. O achado mais significativo dessa parte do estudo, no entanto, não é de magnitude, mas de direção: componentes interharmônicas mostraram-se comparativamente mais recuperáveis do que harmônicas entre os canais mais bem-sucedidos, achado reforçado diretamente pela ablação controlada (Seção~\ref{sec:resultados_ablacao}), na qual o subconjunto formado apenas por interharmônicas atinge o maior $R^2$ de toda a análise de RQ2 ($R^2=0{,}880$) enquanto o subconjunto formado apenas por harmônicas não produz nenhum canal com $R^2$ positivo -- um resultado nunca testado pelos três trabalhos de NILM harmônico revisados nesta pesquisa \citep{djordjevic2017,bouhouras2017,li2019moea} e que confirma, em um contexto de desagregação de cargas real, a hipótese original de \citet{efe2016} sobre a relevância diagnóstica das interharmônicas.}

Tomados em conjunto, RQ1 e RQ2 convergem para uma conclusão mais ampla que motivou esta pesquisa desde a lacuna identificada na Seção~\ref{sec:sintese} do Capítulo~\ref{cap:fundamentacao}: o conteúdo interharmônico, sistematicamente omitido tanto da literatura de NILM baseada em corrente harmônica quanto da literatura clássica de MCSA, carrega informação diagnosticamente relevante tanto para a condição mecânica de máquinas rotativas quanto para a composição de cargas elétricas coexistentes -- e essa relevância se sustenta, nesta pesquisa, em duas tarefas, dois protocolos de avaliação e dois métodos de importância de atributos metodologicamente independentes. A contribuição central desta tese é, portanto, menos a de qualquer resultado numérico isolado do que a de demonstrar, com evidência convergente, que diagnóstico de falha e desagregação de cargas podem ser tratados como duas saídas de um mesmo espaço de entrada espectral agregado, ampliando o escopo de aplicação de uma única cadeia de instrumentação e processamento não intrusivos.

As limitações discutidas na Seção~\ref{cap:limitacoes} -- notavelmente a ausência de defeitos de rolamento na base de dados, o tamanho amostral reduzido de RQ2 e a caracterização apenas qualitativa da severidade de desbalanceamento mecânico -- delimitam o escopo dessas conclusões e apontam diretamente o roteiro de trabalhos futuros mais relevante: a ampliação da base de dados para incluir defeitos de rolamento e mais combinações de carga, permitindo tanto uma cobertura mais completa dos modos de falha mecânica quanto o emprego de arquiteturas sequenciais para RQ2 em paridade com RQ1; a reformulação ordinal ou de regressão do grau de desbalanceamento mecânico, testando diretamente a hipótese levantada na Seção~\ref{sec:discussao_rq1} de que a formulação de classificação discreta subestima a informação diagnóstica disponível para falhas de baixa severidade; e a replicação independente da metodologia em bancadas e instalações distintas das utilizadas nesta pesquisa, condição necessária para avaliar a generalização dos resultados aqui reportados a sistemas trifásicos industriais reais.
\section*{Acknowledgments}
The authors acknowledge financial support from the Fundação para a Ciência e a Tecnologia (FCT, Portugal), I.P., under project UIDB/00048/2020 (DOI: 10.54499/UIDB/00048/2020), and from the National Council for Scientific and Technological Development (CNPq, Brazil) through a Research Productivity Fellowship (Grant No. 301644/2022-5). The authors also thank the Brazilian Federal Agency for Support and Evaluation of Graduate Education (CAPES, Brazil) for postdoctoral fellowship support (Grant Nos. 88887.985910/2024-00 and 88887.985903/2024-00). The authors acknowledge LaMCAD/UFG for providing computational resources, the Instituto Israelita de Ensino e Pesquisa Albert Einstein (IIEP, São Paulo, Brazil) for providing imaging data and research access, and the Research Support Foundation of the State of Goiás (FAPEG, Brazil) for supporting this work through a master's scholarship (Grant No. 202310267000733).
%==============================================================
\section*{Compliance with ethical standards}
%% [REVISAR] Este parágrafo e o de "Data and code availability" abaixo mencionam um
%% comitê de ética (IIEP) e uma base de dados de cintilografia miocárdica (PhysioNet)
%% que não têm relação aparente com este artigo (bancadas de motor/retificador,
%% sem dados de pacientes/imagens médicas). Também o "CRediT authorship contribution
%% statement" lista iniciais de autores (J.C.S.O., G.C.C.N., F.B.F.C. etc.) que não
%% correspondem aos autores declarados no início do documento (Silva, Gomes Pacheco,
%% Rodrigues, Coimbra, Calixto). Tudo isso tem cara de texto herdado do exemplo do
%% template MEDIMA (a classe é originalmente da revista Medical Image Analysis) e
%% não foi adaptado para este trabalho. Não alterei o conteúdo porque não tenho como
%% saber os dados corretos (protocolo de ética real, se houver, dataset real, e a
%% autoria real de cada contribuição) -- mas revise com atenção antes de submeter,
%% para não constar informação de outro artigo/projeto na versão final.
This study was approved by the Research Ethics Committee of the Instituto Israelita de Ensino e Pesquisa Albert Einstein (IIEP), under the CAAE protocol CAAE 65628422.7.0000.0071. All procedures were carried out according to the ethical standards of the Declaration of Helsinki of 1964 and its subsequent amendments.
%==============================================================
\section*{Data and code availability}
The code used to implement the training and inference procedures, as well as the preprocessing scripts, will be made publicly available upon acceptance. The annotated dataset used in this study has been pre-approved for publication and will be hosted on PhysioNet under the following citation:
Nogueira, S., Luz, F., Ortiz de Camargo, T. F., \& Calixto, W. (2025). \textit{Myocardial perfusion scintigraphy image database} (version 1.0.1). PhysioNet. RRID:SCR\_007345. \url{https://doi.org/10.13026/*****}.
These resources will ensure full reproducibility of the experiments, in accordance with the data sharing policies of \textit{Measurement}.
%==============================================================
\section*{Declaration of competing interest}
The authors declare that they have no known competing financial interests or personal relationships that could have appeared to influence the work reported in this article.
%==============================================================
\section*{Funding}
This work was supported by the Fundação para a Ciência e a Tecnologia (FCT, Portugal) under the project UIDB/00048/2020. Additional support was provided by the National Council for Scientific and Technological Development (CNPq, Brazil) through a Research Productivity Fellowship (Grant No. 301644/2022-5), by the Brazilian Federal Agency for Support and Evaluation of Graduate Education (CAPES, Brazil) through postdoctoral fellowships (Grant Nos. 88887.985910/2024-00 and 88887.985903/2024-00), and by the Research Support Foundation of the State of Goiás (FAPEG, Brazil) through a master's scholarship (Grant No. 202310267000733).
%==============================================================
\section*{CRediT authorship contribution statement}
J.C.S.O., G.C.C.N., F.B.F.C.: Validation.
F.A.B.L., T.F.O.C., V.M.G.P., M.R.C.R.: Conceptualization, Methodology, Software, Validation, Formal analysis, Investigation, Resources, Data curation, Writing — original draft, Writing — review \& editing, Visualization.
S.A.N., C.G.R., M.S.R.M.: Conceptualization, Methodology, Software, Validation, Formal analysis, Investigation, Resources, Data curation, Writing — original draft, Writing — review \& editing, Visualization, Supervision.
A.P.C., W.P.C.: Conceptualization, Methodology, Software, Validation, Formal analysis, Investigation, Resources, Data curation, Writing — original draft, Writing — review \& editing, Visualization, Supervision, Project administration, Funding acquisition.
\section*{Appendix A. Supplementary material}
Supplementary data associated with this article can be found in the online version.
%==============================================================
\textcolor{blue}{
\section*{Apêndice B. Checklist de reprodutibilidade}
Este apêndice consolida, em formato de checklist, os elementos necessários à reprodução dos experimentos descritos nesta tese, atendendo à observação do orientador de que informações de reprodutibilidade estavam dispersas ou ausentes no texto. Cada item aponta para a seção onde o dado é discutido em detalhe, ou é marcado como \textbf{[Pendente]} quando o dado ainda não foi confirmado em nenhuma fonte primária (placa de identificação, manual técnico ou registro experimental) -- itens pendentes não foram estimados ou inventados para preencher esta lista.

\textit{Instrumentação e sensores.} Analisador de qualidade de energia IMS PowerNET~PQ-700~G4 (modelo PQK7-60-6C7J-D1N-0-3), sonda de corrente flexível Flex~1000-e, comunicação USB, 512~MB de memória interna, consumo 20~VA, faixa de tensão 50--300~V\textsubscript{ca} (fase-neutro)/50--500~V\textsubscript{ca} (fase-fase) -- Seção~\ref{sec:bancadas}. Classe de exatidão do analisador: \textbf{[Pendente]}, a confirmar no manual técnico completo ou no certificado de calibração do equipamento.

\textit{Motor (bancada RQ1).} WEG W22~Plus, 0,75~kW (1,0~cv), 1.730~rpm (4 polos, 60~Hz), tensão dupla 220/380~V ($\Delta$/Y), corrente nominal 3,02~A (220~V)/1,75~A (380~V), IP55, regime S1, rolamentos 6204-ZZ/6203-ZZ.EM, acionamento por partida direta -- Seção~\ref{subsec:bancada_rq1}. Massa do motor: \textbf{[Pendente]} (campo em branco na placa de identificação; obter por pesagem direta ou catálogo WEG).

\textit{Método de indução de falha mecânica.} Desbalanceamento induzido por remoção controlada de parafusos de um disco calibrado (rótulo \texttt{RET~N~PARAF} = número de parafusos retirados) -- Seção~\ref{sec:discussao_rq1}. Massa, raio e posição angular de cada parafuso removido: \textbf{[Pendente]} -- já registrado como limitação estrutural na Seção~\ref{cap:limitacoes}, pois impede caracterização quantitativa (não apenas qualitativa) do grau de desbalanceamento. Método de indução de desalinhamento: \textbf{[Pendente]}.

\textit{Bancada de retificador (RQ2).} Topologia exata do retificador trifásico, cargas não lineares utilizadas e critério de composição das combinações de carga (incluindo as condições \texttt{-DIODO~2}, \texttt{-DIODO~3}, \texttt{-DIODO~6} e \texttt{RETIFICADOR~COMPLETO}): \textbf{[Pendente]} -- Seção~\ref{subsec:bancada_rq2}. Se a mesma técnica de multiplicação de corrente por enrolamento triplo empregada na bancada RQ1 (Seção~\ref{subsec:bancada_rq1}) também foi aplicada aqui: \textbf{[Pendente de confirmação]}.

\textit{Protocolo de aquisição.} Intervalo de registro (\textit{datalog}) de 10~s entre aquisições consecutivas, cada uma agregada internamente pelo analisador sobre uma janela de 200~ms conforme IEC~61000-4-7, aquisição simultânea de todos os canais a 16~bits e 10.240 amostras/s, agrupamento espectral Subgrupo (S-Grupo) -- Seções~\ref{sec:bd} e~\ref{sec:atributos}. Base de dados resultante: 19.501 aquisições em 17 condições experimentais (Seção~\ref{sec:resultados_bd}), cada condição correspondendo a uma única sessão de aquisição contínua (sem repetições independentes -- limitação central discutida nas Seções~\ref{sec:protocolo} e~\ref{cap:limitacoes}). Duração de cada sessão: da ordem de horas por condição (variável entre condições; duração exata por rótulo não tabulada neste texto -- \textbf{[Pendente]}). Ordem de execução dos experimentos entre as 17 condições: \textbf{[Pendente]}. Sincronização entre canais de tensão e corrente e entre planilhas de harmônicas e inter-harmônicas: aquisição simultânea nativa do analisador para tensão/corrente; pareamento por \textit{timestamp} mais próximo entre as duas planilhas de origem, com tolerância registrada como \texttt{join\_delta\_seconds} no pipeline (Seção~\ref{sec:atributos}).

\textit{Pré-processamento e engenharia de atributos.} 984 canais espectrais candidatos (666 harmônicos + 318 inter-harmônicos), incluindo atributos derivados de distorção agregada (THD, THDe) -- Seção~\ref{sec:atributos}. Tratamento de valores ausentes/\textit{outliers}: \textbf{[Pendente]} de documentação explícita no texto principal. Escalonamento: \texttt{StandardScaler}, ajustado por \textit{split}/dobra dentro do \textit{pipeline} de cada modelo sensível a escala -- Seção~\ref{sec:protocolo}.

\textit{Seleção de atributos.} F-score ANOVA com eliminação de redundância por correlação (limiar 0,95), reajustada a cada \textit{split}/dobra usando exclusivamente as linhas de treino -- Seção~\ref{sec:protocolo}. Para RQ2, seleção adicionalmente reajustada excluindo as quatro condições combinadas inteiras, não apenas a condição retida por dobra -- mesma seção.

\textit{Modelos e hiperparâmetros.} Seis modelos clássicos (Random~Forest, SVM, XGBoost, MLP, $k$-NN, Regressão~Logística) com hiperparâmetros padrão das implementações \textit{scikit-learn}/\textit{XGBoost}, sem busca sistemática de hiperparâmetros na comparação principal -- Seções~\ref{subsec:modelos_classicos} e~\ref{sec:protocolo}. Busca aleatória de hiperparâmetros demonstrada isoladamente para a LSTM (\texttt{units} $\in\{32,64,128\}$, \textit{dropout} $\in\{0{,}1;0{,}2;0{,}3\}$) -- Seção~\ref{subsec:rnns}. Seis arquiteturas sequenciais (SimpleRNN, LSTM, GRU, BiLSTM, CNN+LSTM, Transformer) especificadas na Tabela~\ref{tab:res_rnn_spec}, janela de 8 aquisições, passo 4, treinadas por até 20 épocas com parada antecipada (paciência 5) -- Seções~\ref{sec:resultados_rnn} e~\ref{sec:protocolo}.

\textit{Sementes aleatórias.} Modelos clássicos: \texttt{random\_state=42}, fixo entre execuções. Arquiteturas sequenciais: originalmente semente única por arquitetura (limitação registrada na Seção~\ref{cap:limitacoes}); retreinadas sob cinco sementes independentes por arquitetura (Seção~\ref{sec:resultados_rnn}) para estimar a variância de desempenho entre inicializações.

\textit{Hardware e versões de software.} \textbf{[Pendente]} -- versões de Python, \textit{scikit-learn}, XGBoost, TensorFlow/Keras, sistema operacional e especificação de hardware (CPU/GPU) utilizados nos experimentos ainda não foram registradas neste texto; a documentar antes da versão final, junto com a disponibilização do código mencionada na seção ``Data and code availability''.
}
%==============================================================
\textcolor{blue}{
\section*{Apêndice C. Síntese em linguagem acessível}
Este apêndice reúne, sem jargão estatístico, uma síntese dos achados centrais de RQ1 e RQ2 discutidos no Capítulo~\ref{cap:resultados} -- deslocada do núcleo do capítulo de Resultados (Seção~\ref{sec:resultados_sensibilidade}) para este apêndice, atendendo à observação do orientador de que essa síntese, embora útil como material de apoio a uma apresentação de qualificação para público não especializado, não deveria competir com a evidência quantitativa primária na leitura do capítulo de Resultados. As quatro figuras a seguir reapresentam dados já discutidos no corpo do texto, com enquadramento simplificado, e não constituem evidência adicional.

\begin{figure}[H]
    \centering
    \includegraphics[width=0.68\textwidth]{leigos_accuracy}
    \caption{Taxa de acerto do diagnóstico automático: 92\% ao misturar defeitos elétricos e mecânicos \textit{vs.} 81\% ao isolar somente defeitos mecânicos, mais sutis.}
    \label{fig:res_leigos_acc}
\end{figure}

\begin{figure}[H]
    \centering
    \includegraphics[width=0.65\textwidth]{leigos_electrical}
    \caption{Assinatura de corrente (H2--H15): motor saudável \textit{vs.} motor com defeito elétrico proposital.}
    \label{fig:res_leigos_elec}
\end{figure}

\begin{figure}[H]
    \centering
    \includegraphics[width=0.68\textwidth]{leigos_voltage_share}
    \caption{Fração de tensão entre os 15 canais mais importantes para o diagnóstico: 1 em 15 com defeitos elétricos e mecânicos misturados, contra 5 em 15 com apenas defeitos mecânicos.}
    \label{fig:res_leigos_volt}
\end{figure}

\begin{figure}[H]
    \centering
    \includegraphics[width=0.68\textwidth]{leigos_nilm}
    \caption{Qualidade de separação (NILM) para os cinco canais espectrais em que a desagregação funcionou melhor.}
    \label{fig:res_leigos_nilm}
\end{figure}
}
%==============================================================
% \bibliographystyle{elsarticle-harv}
%\bibliographystyle{model1a-num-names.bst}
%\biboptions{authoryear}
\bibliographystyle{elsarticle-num-names}
\bibliography{01_publicacao.bib}
\end{document}
